\SourcePage{280}
\SGAsection{4}{由子群 ${}_nT$ 刻画子环面 $T$}
\textbf{1. 主要\Term{theorem}{定理}的陈述}
\begin{SGAresult}{\Term{theorem}{定理}}{4.1}
设 $S$ 是\Term{locally-noetherian}{局部 noether}的\Term{connected}{连通}\Term{prescheme}{预概形}，$G$ 是 $S$ 上\Term{finite-type}{有限型}的 $S$-\Term{group-prescheme}{群预概形}，$q$ 是在 $S$ 上\Term{invertible}{可逆}的整数，且 $q>1$，$r$ 是非负整数。对于每个为 $q$ 的幂的整数 $n$，设 $M(n)$ 是 $G$ 的一个\Term{multiplicative-type}{乘法型}\Term{group-subscheme}{子群概形}，其型为 $(\mathbb Z/n\mathbb Z)^r$。假设：

a) \Term{subgroup}{子群}族 $M(n)$ 相容，即若整数 $m$ 整除 $n$，则
\[{}_mM(n)=M(m).\]
b) 存在 $S$ 的一点 $s$ 及 $G_s$ 的一个\Term{subtorus}{子环面} $T_s$，使
\[M(n)_s={}_nT_s\qquad\text{对所有 }n.\]
c) 对于 $S$ 的每一点 $t$，存在 $G_t$ 的一个\Term{closed-subscheme}{闭子概形} $F_{\bar t}$，它\Term{affine}{仿射}，并对所有 $n$ 都包含 $M(n)_{\bar t}$。

则存在唯一的 $G$ 的\Term{subtorus}{子环面} $T$，使得对每个为 $q$ 的幂的整数 $n$，都有 ${}_nT=M(n)$。
\end{SGAresult}
关于\Term{morphism}{态射}的\Term{lifting}{提升}，也有一个类似\Term{theorem}{定理}：
\begin{SGAresult}{\Term{theorem}{定理}}{4.1 bis}
设 $S,G,q$ 如上，$T$ 是一个 $S$-\Term{torus}{环面}；对于每个为 $q$ 的幂的整数 $n$，设 $u(n):{}_nT\to G$ 是一个 $S$-\Term{group-morphism}{群态射}。假设：

a) \Term{morphism}{态射}族 $u(n)$ 相容，即若 $m$ 整除 $n$，则
\[u(m)=u(n)|_{{}_mT}.\]
b) 存在 $S$ 的一点 $s$ 及\Term{group-morphism}{群态射}
\[u_s:T_s\longrightarrow G_s,\]
使得对每个为 $q$ 的幂的整数 $n$，都有 $u_s|_{{}_nT_s}=u(n)_s$。
\SourcePage{281}
c) 对于 $S$ 的每一点 $t$，存在 $G_{\bar t}$ 的一个\Term{closed-subscheme}{闭子概形} $F_{\bar t}$，它\Term{affine}{仿射}，并对所有 $n$ 都包含 $u(n)_{\bar t}({}_nT_{\bar t})$。

则存在唯一的\Term{group-morphism}{群态射} $u:T\to G$，使得对每个为 $q$ 的幂的整数 $n$，$u$ 在 ${}_nT$ 上的\Term{restriction}{限制}等于 $u(n)$。
\end{SGAresult}
\begin{SGAresult}{注记}{4.2}
利用\Term{trait}{离散赋值环谱}上\Term{flat}{平坦}\Term{finite-type}{有限型}\Term{group-prescheme}{群预概形}的\Term{abelian-rank}{阿贝尔秩}的\Term{lower-semicontinuous}{下半连续}性（参见第 X 讲 8.7），在 4.1 及 4.1 bis 的陈述中，可以减弱条件 c)，只要求对 $S$ 的每个\Term{maximal-point}{极大点} $t$，存在所要求的\Term{affine}{仿射}\Term{closed-subscheme}{闭子概形} $F_{\bar t}$。
\end{SGAresult}
下面说明如何由 4.1 \Term{pushout}{推出} 4.1 bis。令 $G'=G\times_ST$。对于每个为 $q$ 的幂的整数 $n$，考察\Term{group-morphism}{群态射}
\[v(n):{}_nT\longrightarrow G',\]
它到 $G$ 及 $T$ 的\Term{projection}{投影}分别为 $u(n)$ 及\Term{canonical}{典范}\Term{immersion}{浸入} ${}_nT\to T$。因此，\Term{morphism}{态射} $v(n)$ 是\Term{immersion}{浸入}，以 $M(n)$ 表示其像\Term{subgroup}{子群}。显然，\Term{subgroup}{子群}族 $M(n)$ 在 4.1 的意义下相容；群 $M(n)_s$ 等于 ${}_nT'_s$，其中 $T'_s$ 是 $G'_s$ 中作为 $u_s$ 的图的\Term{subtorus}{子环面}；对于 $S$ 的每一点 $t$，$G'_{\bar t}$ 的\Term{affine}{仿射}\Term{closed-subscheme}{闭子概形} $F_{\bar t}\times_{\bar t}T_{\bar t}$ 对所有 $n$ 都包含\Term{subgroup}{子群} $M(n)_{\bar t}$。根据 4.1，存在 $G'$ 的一个\Term{subtorus}{子环面} $T'$，使得对每个为 $q$ 的幂的整数 $n$，都有 ${}_nT'=M(n)$。以 $f$ 表示从 $G'$ 到 $T$ 的\Term{projection}{投影}在 $T'$ 上的\Term{restriction}{限制}，$f(n)$ 表示 $f$ 在 $M(n)$ 上的\Term{restriction}{限制}。则
\[f(n)\circ v(n)=\operatorname{id}_{{}_nT}.\]
$T'$ 在 $s$ 处的\Term{fiber}{纤维}就是已经记为 $T'_s$、等于 $u_s$ 的图的\Term{torus}{环面}（这由第 IX 讲 4.8 b) 得到），故 $f_s$ 为\Term{isomorphism}{同构}。但是，$\operatorname{Ker}f$ 及 $\operatorname{Coker}f$ 都是\Term{multiplicative-type}{乘法型}群（第 IX 讲 2.7），且因 $S$ \Term{connected}{连通}而具有\Term{constant-type}{常型}，因而均为\Term{unit-group}{单位群}，$f$ 是\Term{isomorphism}{同构}。令 $v$ 为 $f$ 的逆\Term{isomorphism}{同构}，则
\[v|_{{}_nT}=v(n).\]
因此，$v$ 与从 $G'$ 到 $G$ 的\Term{projection}{投影}的复合是一个满足要求的\Term{morphism}{态射} $u:T\to G$。以上证明了 $u$ 的存在性；其唯一性则无论如何都由第 IX 讲 4.8 a) 得到。

\textbf{2. 应用}

我们的目标是推广第 IX 讲的\Term{theorem}{定理} 7.1。

设 $A$ 是\Term{complete}{完备} noether \Term{local-ring}{局部环}，$I$ 为其\Term{maximal-ideal}{极大理想}，$S=\operatorname{Spec}A$，$S_m=\operatorname{Spec}(A/I^m)$。对于每个\Term{prescheme}{预概形} $X$，置 $X_m=X\times_SS_m$。

再设 $G$ 是\Term{finite-type}{有限型} $S$-\Term{group-prescheme}{群预概形}，$T$ 是 $S$-\Term{torus}{环面}，$q$ 是在 $S$ 上\Term{invertible}{可逆}的整数，$u_m:T_m\to G_m$（$m\in\mathbb N$）是一族相容的\Term{group-morphism}{群态射}。当整数 $n$ 遍历 $q$ 的幂时，以 $u_m(n)$ 表示 $u_m$ 在 ${}_nT_m$ 上的\Term{restriction}{限制}，$u(n)$ 表示唯一的\Term{group-morphism}{群态射}
\[u(n):{}_nT\longrightarrow G,\]
\SourcePage{282}
它对所有 $m$ 都延拓 $u_m(n)$（1.6 a)）。如果对 $S$ 的每一点 $t$，存在 $G_{\bar t}$ 的一个\Term{affine}{仿射}\Term{closed-subprescheme}{闭子预概形} $F_{\bar t}$，使它对于每个为 $q$ 的幂的整数 $n$ 都包含 $u(n)_{\bar t}({}_nT_{\bar t})$，我们就称族 $(u_m)_{m\in\mathbb N}$ 是\Term{admissible-family}{容许的}（下面会看到，这个性质与 $q$ 的选择无关）。
\begin{SGAresult}{命题}{4.3}
采用以上记号，\Term{canonical}{典范}映射
\[\operatorname{Hom}_{S\text{-gr}}(T,G)\longrightarrow\varprojlim_m\operatorname{Hom}_{S_m\text{-gr}}(T_m,G_m)\]
给出左端与右端中由\Term{admissible-family}{容许的}相容族组成的\Term{subset}{子集}之间的\Term{isomorphism}{同构}。
\end{SGAresult}
事实上，只须应用 4.1 bis，取 $s$ 为 $S$ 的\Term{closed-point}{闭点}。
\begin{SGAresult}{推论}{4.4}
采用以上记号，进一步假设 $G$ 的\Term{fiber}{纤维}均\Term{affine}{仿射}，则\Term{canonical}{典范}映射
\[\operatorname{Hom}_{S\text{-gr}}(T,G)\longrightarrow\varprojlim_m\operatorname{Hom}_{S_m\text{-gr}}(T_m,G_m)\]
是\Term{isomorphism}{同构}。
\end{SGAresult}
事实上，如果 $G$ 的\Term{fiber}{纤维}均\Term{affine}{仿射}，则任意相容族都是\Term{admissible-family}{容许的}。
\begin{SGAresult}{注记}{4.5}
当 $G$ \Term{separated}{分离}时，在 4.3 中可以把环 $A$ 替换为一个 $I$-进 noether 环；此时可以使用 EGA III 5.4.1 代替 1.6 a)。
\end{SGAresult}
\textbf{3. 4.1 的证明}
\begin{SGAresult}{\Term{lemma}{引理}}{4.6}
设 $k$ 是域，$G$ 是 $k$-\Term{algebraic-group}{代数群}，$q$ 是与 $k$ 的\Term{residue-characteristic}{剩余特征}互素的整数，$r>0$ 是整数，$M(n)$（$n$ 遍历 $q$ 的幂）是 $G$ 的一族相容的\Term{multiplicative-type}{乘法型}\Term{subgroup}{子群}，其型为 $(\mathbb Z/n\mathbb Z)^r$。则存在 $G$ 的最小\Term{closed-subscheme}{闭子概形} $H$，对所有 $n$ 都包含 $M(n)$。此外，$H$ 是 $G$ 的一个\Term{smooth}{光滑}、\Term{connected}{连通}且交换的\Term{algebraic-subgroup}{代数子群}，并且“其构造与\Term{base-field-extension}{基域扩张}相容”。
\end{SGAresult}
\begin{proof}
令 $M$ 为 $\operatorname{ens}(G)$ 的\Term{subset}{子集}，由所有\Term{subgroup}{子群} $M(n)$ 的\Term{underlying-set}{底集}合之并组成；令 $H$ 为 $G$ 中以 $M$ 的\Term{closure}{闭包}为底\Term{space}{空间}的\Term{reduced}{约化}\Term{closed-subscheme}{闭子概形}。由于 $M(n)$ 在 $k$ 上\Term{etale}{平展}，从而\Term{reduced}{约化}，$M(n)$ 包含于 $H$，所以 $H$ 是 $G$ 中对所有 $n$ 都包含 $M(n)$ 的最小\Term{closed-subscheme}{闭子概形}。现取 $k$ 的一个\Term{algebraically-closed}{代数闭}\Term{group-extension}{扩张} $\bar k$。按构造，\Term{subscheme}{子概形}族 $M(n)$ 在 $H$ 中\Term{schematically-dense}{概形稠密}（第 IX 讲 4.1），故 $M(n)_{\bar k}$ 在 $H_{\bar k}$ 中\Term{schematically-dense}{概形稠密}（第 IX 讲 4.5）；另一方面，$M(n)_{\bar k}$ \Term{reduced}{约化}，从而 $H_{\bar k}$ 必然等于 $G_{\bar k}$ 中以 $M_{\bar k}$ 的\Term{closure}{闭包}为底\Term{space}{空间}的\Term{reduced}{约化}\Term{closed-subscheme}{闭子概形}。这证明 $H$ \Term{geometrically-reduced}{几何约化}（EGA IV 4.6.1）。由于族 $M(n)$ 相容，$M$ 对\Term{group-law}{群运算}封闭；又由于 $M\times_kM$ 在 $\operatorname{ens}(H\times_kH)$ 中\Term{dense}{稠密}，且 $H\times_kH$ \Term{reduced}{约化}（EGA IV 4.6.5），立即\Term{pushout}{推出} $H$ 是 $G$ 的一个\Term{algebraic-subgroup}{代数子群}。此外，$H$ 因\Term{geometrically-reduced}{几何约化}而在 $S$ 上\Term{smooth}{光滑}（第 VI$_A$ 讲 1.3.1），又因各 $M(n)$ 交换而为\Term{abelian-group}{交换群}。还须证明 $H$ \Term{connected}{连通}。
\SourcePage{283}
令 $H^0$ 为 $H$ 的\Term{identity-component}{单位连通分支}，$m$ 为 $H/H^0$ 的\Term{geometric-point}{几何点}个数，$q^s$ 为 $m$ 的素因子\Term{resolution}{分解}中 $q$ 的幂。对于每个为 $q$ 的幂的整数 $n$，$q^sM(n)$ 都包含在 $H^0$ 中。但族 $M(n)$ 相容，且 $M(n)$ 的型为 $(\mathbb Z/n\mathbb Z)^r$，所以 $q^sM(nq^s)=M(n)$。因此，$H^0$ 对所有 $n$ 都包含 $M(n)$，从而等于 $H$。
\end{proof}
以上结果使我们可以进一步说明 4.1 的条件 c)。事实上，根据上面的论证，可以考察 $G_t$ 中对所有 $n$ 都包含 $M(n)_t$ 的最小\Term{closed-subscheme}{闭子概形} $H_t$。$H_t$ 的构造与\Term{base-field-extension}{基域扩张}交换（4.6），所以 $H_{\bar t}=H_t\times_t\bar t$ 包含于\Term{affine}{仿射}闭集 $F_{\bar t}$，因而\Term{affine}{仿射}。总之，$H_t$ \Term{affine}{仿射}。另一方面，我们知道（4.6），$H_{\bar t}$ 是\Term{smooth}{光滑}、\Term{connected}{连通}的交换\Term{algebraic-group}{代数群}。由\Term{algebraically-closed-field}{代数闭域}上\Term{affine}{仿射}、\Term{smooth}{光滑}、\Term{connected}{连通}的交换\Term{algebraic-group}{代数群}的结构（BIBLE，4，\Term{theorem}{定理} 4），$H_{\bar t}$ 是一个\Term{torus}{环面} $T_{\bar t}$ 与一个\Term{unipotent-group}{幺幂群} $U_{\bar t}$ 的\Term{direct-product}{直积}。但 $M(n)_{\bar t}$ 必然包含于 $T_{\bar t}$，因此 $H_{\bar t}=T_{\bar t}$，从而 $H_t$ 是\Term{torus}{环面}。

现在可以着手证明 4.1。

\textbf{a) 解的唯一性。}只须应用第 IX 讲 4.8 b)。

\textbf{b) $S$ 为\Term{residue-field}{剩余域} $k$ \Term{algebraically-closed}{代数闭}的\Term{complete-discrete-valuation-ring}{完备离散赋值环} $A$ 的谱的情形。}以 $K$ 表示 $A$ 的\Term{field-of-fractions}{分式域}，$s$ 表示 $S$ 的\Term{closed-point}{闭点}，$t$ 表示其\Term{generic-point}{泛点}，$J$ 表示 $A$ 的\Term{maximal-ideal}{极大理想}，$S_m=\operatorname{Spec}(A/J^m)$，$X_m=X\times_SS_m$。

分两种情形：

\emph{第一种情形：}4.1 b) 中的点 $s$ 是 $S$ 的\Term{generic-point}{泛点} $t$。令 $T'$ 为\Term{torus}{环面} $T_t$ 在 $G$ 中的\Term{schematic-closure}{概形闭包}。于是，闭\Term{fiber}{纤维} $T'_s$ 是 $G_s$ 的一个 $r$ 维\Term{algebraic-subgroup}{代数子群}，它对所有 $n$ 都包含 $M(n)_s$，故 $T'_s$ 包含 $H_s$。但 $H_s$ 是一个包含 $M(n)_s$ 的\Term{torus}{环面}，因而 $H_s$ 的秩至少为 $r$。所以，$H_s$ 与 $T'_s$ 的\Term{identity-component}{单位连通分支}具有相同底\Term{space}{空间}。于是 $T'$ 的“\Term{identity-component}{单位连通分支}” $(T')^0$ 是 $G$ 的一个\Term{group-subprescheme}{子群预概形}，它\Term{flat}{平坦}、\Term{separated}{分离}（第 VI$_B$ 讲 5.2），其泛\Term{fiber}{纤维} $T_t$ 是\Term{torus}{环面}，\Term{reduced}{约化}的闭\Term{fiber}{纤维} $H_s$ 也是\Term{torus}{环面}。于是 $(T')^0$ 是\Term{torus}{环面}（第 X 讲 8.8），记为 $T$。群 ${}_nT$ 及 $M_n$ 都在 $S$ 上\Term{smooth}{光滑}，从而\Term{reduced}{约化}；它们具有相同底\Term{space}{空间}，故相等。因此，\Term{torus}{环面} $T$ 就是问题的解。

\emph{第二种情形：}4.1 b) 中的点 $s$ 是 $S$ 的\Term{closed-point}{闭点} $s$。用包含族 $M(n)_t$ 的最小\Term{algebraic-subgroup}{代数子群} $H_t$（4.6）在 $G$ 中的\Term{schematic-closure}{概形闭包}替换 $G$，可以假定 $G_t$ \Term{affine}{仿射}。

对于每个整数 $m\geq0$，由 2.2，存在 $G_m$ 的唯一一个\Term{subtorus}{子环面} $T_m$ \Term{lifting}{提升} $T_s$，并且对每个为 $q$ 的幂的整数 $n$，都有 ${}_nT_m=M(n)_m$。另一方面，令 $T'=\mathbb G_{m,S}^r$。由于 $k$ \Term{algebraically-closed}{代数闭}，$T_s$ \Term{split-torus}{平凡}，故存在 $k$-\Term{isomorphism}{同构} $u_s:T'_s\xrightarrow{\sim}T_s$。\Term{morphism}{态射} $u_s$ 唯一地\Term{lifting}{提升}为 $S_m$-\Term{isomorphism}{同构} $u_m:T'_m\to T_m$（第 IX 讲 3.3）。\Term{morphism}{态射}族 $(u_m)_{m\in\mathbb N}$ 在\Term{limit}{极限}处定义一个\Term{morphism}{态射} $\widehat u$，其源和靶是 $T'$ 及 $G$ 沿闭\Term{fiber}{纤维}的\Term{formal-completion}{形式完备化}
\SourcePage{284}
$\widehat{T'}$ 及 $\widehat G$：
\[
\begin{tikzcd}
\widehat{T'}\arrow[r,"\widehat u"]\arrow[d,"i"']&\widehat G\arrow[d,"j"]\\
T'&G,
\end{tikzcd}
\]
其中 $i,j$ 为\Term{canonical}{典范}\Term{morphism}{态射}。

我们将证明\Term{morphism}{态射} $\widehat u$ 可\Term{algebraization}{代数化}。为此，先归结到群 $G$ \Term{affine}{仿射}的情形。
\begin{SGAresult}{\Term{lemma}{引理}}{4.7}
设 $S$ 是\Term{discrete-valuation-ring}{离散赋值环} $A$ 的谱，$X,Y$ 是两个在 $S$ 上\Term{quasi-compact}{拟紧}、\Term{quasi-separated}{拟分离}且\Term{flat}{平坦}的 $S$-\Term{prescheme}{预概形}。则\Term{canonical}{典范}映射
\[\Gamma(X,\mathcal O_X)\otimes_A\Gamma(Y,\mathcal O_Y)\longrightarrow\Gamma(X\times_SY,\mathcal O_{X\times_SY})\]
是\Term{isomorphism}{同构}。
\end{SGAresult}
\begin{proof}
令 $f:X\to S$（相应地 $g:Y\to S$）为\Term{structural-morphism}{结构态射}。由于 $X$（相应地 $Y$）\Term{quasi-compact}{拟紧}且\Term{quasi-separated}{拟分离}，由 EGA I 9.2.2 及 EGA IV 1.7.4，$f_*(\mathcal O_X)$（相应地 $g_*(\mathcal O_Y)$）是一个\Term{quasi-coherent}{拟凝聚} $\mathcal O_S$-\Term{algebra}{代数}，故对应于一个\Term{affine}{仿射} $S$-\Term{scheme}{概形} $\widetilde X$（相应地 $\widetilde Y$）。根据假设，$Y$ 在 $S$ 上\Term{flat}{平坦}；由 EGA III 1.4.15，结合 EGA IV 1.7.21，\Term{canonical}{典范}映射（由\Term{canonical}{典范}\Term{morphism}{态射} $X\to\widetilde X$ \Term{derive}{导出}）
\[\Gamma(\widetilde X\times_SY,\mathcal O_{\widetilde X\times_SY})\longrightarrow\Gamma(X\times_SY,\mathcal O_{X\times_SY})\]
是\Term{isomorphism}{同构}。但是，由于 $X$ 在 $S$ 上\Term{flat}{平坦}，$\widetilde X$ 也在 $S$ 上\Term{flat}{平坦}（因为在 $A$ 上\Term{flat}{平坦}\Term{equivalence}{等价}于\Term{torsion-free}{无挠}）。交换 $X,Y$ 的地位，再应用 EGA III 1.4.15，得到\Term{isomorphism}{同构}
\begin{align*}
\Gamma(\widetilde X,\mathcal O_{\widetilde X})\otimes_A\Gamma(\widetilde Y,\mathcal O_{\widetilde Y})
&\simeq\Gamma(\widetilde X\times_S\widetilde Y,\mathcal O_{\widetilde X\times_S\widetilde Y})\\
&\xrightarrow{\sim}\Gamma(\widetilde X\times_SY,\mathcal O_{\widetilde X\times_SY}),
\end{align*}
\Term{lemma}{引理}得证。
\end{proof}
在当前情形中，$S$-群 $G$ 在 $S$ 上\Term{flat}{平坦}且\Term{finite-type}{有限型}，所以\Term{quasi-compact}{拟紧}且\Term{quasi-separated}{拟分离}。因此可以把\Term{lemma}{引理}应用于 $G\times_SG$：
\[\Gamma(G,\mathcal O_G)\otimes_A\Gamma(G,\mathcal O_G)\xrightarrow{\sim}\Gamma(G\times_SG,\mathcal O_{G\times_SG}).\]
于是，定义 $G$ 中乘法的\Term{morphism}{态射} $m_G:G\times_SG\to G$ 对应于一个\Term{morphism}{态射}
\[\Gamma(G,\mathcal O_G)\longrightarrow\Gamma(G\times_SG,\mathcal O_{G\times_SG})\xrightarrow{\sim}\Gamma(G,\mathcal O_G)\otimes_A\Gamma(G,\mathcal O_G),\]
从而对应于一个 $S$-\Term{morphism}{态射}
\[m_{\widetilde G}:\widetilde G\times_S\widetilde G\longrightarrow\widetilde G,\]
其中 $\widetilde G$ 表示以 $\Gamma(G,\mathcal O_G)$ 为 $A$-\Term{algebra}{代数}的\Term{affine}{仿射} $S$-\Term{scheme}{概形}。由此\Term{form}{形式}地检查可知，$m_{\widetilde G}$ 赋予 $\widetilde G$ 一个 $S$-\Term{group-scheme}{群概形}结构，使得\Term{canonical}{典范}\Term{morphism}{态射} $v:G\to\widetilde G$ 是 $S$-\Term{group-morphism}{群态射}。
\SourcePage{285}
\begin{SGAresult}{注记}{4.8}
$\widetilde G$ 起到 $G$ 的最大“\Term{affine-quotient}{仿射商}”的\Term{action}{作用}；而且可以证明，$\widetilde G$ 确实是 $G$ 对 fpqc 的商，因而在 $S$ 上\Term{finite-type}{有限型}（第 XVII 讲附录 III，2）。
\end{SGAresult}
在当前情形中，$G$ 的泛\Term{fiber}{纤维} $G_t$ \Term{affine}{仿射}，故由 EGA I 9.3.3，$v_t$ 是\Term{isomorphism}{同构}。由于 $\widetilde G$ \Term{affine}{仿射}，相容\Term{morphism}{态射}族
\[w_m=v_mu_m:T'_m\longrightarrow\widetilde G_m\qquad(m\in\mathbb N)\]
来自唯一的 $S$-\Term{group-morphism}{群态射}（第 IX 讲 7.1）
\[w:T'\longrightarrow\widetilde G.\]
令 $T_t$ 为 $G_t$ 中等于 $v_t^{-1}w_t(T'_t)$ 的\Term{subtorus}{子环面}。作为一个秩为 $r$ 的\Term{torus}{环面}的像，\Term{torus}{环面} $T_t$ 的秩至多为 $r$。下面证明 $T_t$ 对所有 $n$ 都包含 $M(n)_t$。事实上，令 $u(n)_m$ 为把 $u_m$ \Term{restriction}{限制}到 $({}_nT')_m$ 得到的 $S_m$-\Term{isomorphism}{同构}
\[({}_nT')_m\xrightarrow{\sim}M(n)_m.\]
相容\Term{morphism}{态射}族 $u(n)_m$ 来自唯一的 $S$-\Term{isomorphism}{同构} $u(n):{}_nT'\xrightarrow{\sim}M(n)$（因为 $M(n)$ 在 $S$ 上\Term{finite}{有限}）。对每个整数 $m\geq0$，于是有等式
\[w_m|_{({}_nT')_m}=(v_mu_m)|_{({}_nT')_m}=(v\circ u(n))_m.\]
因此，$w|_{{}_nT'}=v\circ u(n)$（1.6 a)）。特别地，
\[w_t|_{({}_nT')_t}=v_t\circ u(n)_t,\]
故 $v_t^{-1}w_t({}_nT')_t=u(n)_t({}_nT')_t=M(n)_t$。

这就证明 $T_t$ 包含 $M(n)_t$，并\Term{pushout}{推出} $T_t$ 的秩为 $r$。如已处理的第一种情形那样，考察 $T_t$ 在 $G$ 中的\Term{schematic-closure}{概形闭包} $T''$，即可完成证明。由于 $T_t$ 包含 $M(n)_t$，$T''_s$ 包含 $M(n)_s$，因而包含 $T_s$（\Term{dense}{稠密}性\Term{theorem}{定理}）。另一方面，由于 $T''$ 在 $S$ 上\Term{flat}{平坦}且 $T_t$ 为 $r$ 维，$T''_s$ 为 $r$ 维（第 VI$_B$ 讲 4）。总之，$T_s$ 与 $T''_s$ 的\Term{identity-component}{单位连通分支}具有相同底\Term{space}{空间}；如第一种情形那样，可知 $T''$ 的\Term{identity-component}{单位连通分支}是满足要求的 $G$ 的\Term{subtorus}{子环面}。

\textbf{c) 完成 4.1 的证明。}

为证明\Term{subtorus}{子环面} $T$ 存在，可以假定 $S$ \Term{reduced}{约化}（2.2）。由 3.1 (ii) $\Rightarrow$ (i)，只须证明：$S$ 中那些满足“存在 $G_x$ 的\Term{subtorus}{子环面} $T_x$，使对每个为 $q$ 的幂的整数 $n$ 都有 ${}_nT_x=M(n)_x$”的点 $x$ 所成的\Term{set}{集合} $U$ 等于 $\operatorname{ens}(S)$。这里的\Term{torus}{环面} $T_x$ 必然唯一；由 fpqc \Term{descent}{下降}，只须在\Term{group-extension}{扩张} $x$ 的\Term{residue-field}{剩余域}后证明它存在（第 IX 讲 4.8 b)）。

由于 $S$ \Term{locally-noetherian}{局部 noether}且\Term{connected}{连通}，又由于 $U$ \Term{nonempty}{非空}（它包含 4.1 b) 中的点 $s$），经一个直接的论证，只须证明：若 $x,x'$ 是 $S$ 的两点，$x$ 为 $x'$ 的\Term{specialization}{特殊化}，且其中一点属于 $U$，则两点都属于 $U$。用通常的技巧（EGA II 7.1.9），归结到 $S$ 是一个\Term{residue-field}{剩余域}\Term{algebraically-closed}{代数闭}的\Term{complete-discrete-valuation-ring}{完备离散赋值环}的谱的情形，此情形已在 b) 中处理。
\SourcePage{286}
\Term{theorem}{定理} 4.1 的证明至此完成。
\SGAsection{5}{等于自身连通正规化子的光滑子群函子的可表性}
\begin{SGAresult}{命题}{5.1}
设 $S$ 是\Term{prescheme}{预概形}，$G$ 是\Term{finite-presentation}{有限呈示}的 $S$-\Term{group-prescheme}{群预概形}，$H$ 是 $G$ 的一个\Term{group-subprescheme}{子群预概形}，它\Term{smooth}{光滑}且\Term{fiber}{纤维}\Term{connected}{连通}。则：

a) $H$ 在 $G$ 中的\Term{normalizer}{正规化子} $N$ 可由 $G$ 的一个在 $S$ 上\Term{finite-presentation}{有限呈示}的\Term{closed-subprescheme}{闭子预概形}表示。

b) 下列条件\Term{equivalence}{等价}：

i) \Term{canonical}{典范}\Term{immersion}{浸入} $H\to N$ 是\Term{open-immersion}{开浸入}。

ii) 群 $N$ 沿\Term{unit-section}{单位截面}\Term{smooth}{光滑}，其\Term{identity-component}{单位连通分支}（此时\Term{representability}{可表}，第 VI$_B$ 讲 3.10）等于 $H$。

iii) 对 $S$ 的每一点 $s$，都有 $H_s=(N_s)^0$。
\end{SGAresult}
\begin{proof}
\Term{group-prescheme}{群预概形} $H$ 在 $S$ 上\Term{locally-finite-presentation}{局部有限呈示}（因为 $H$ 在 $S$ 上\Term{smooth}{光滑}），且\Term{fiber}{纤维}\Term{connected}{连通}，故在 $S$ 上\Term{finite-presentation}{有限呈示}（第 VI$_B$ 讲 5.3.3）。断言 a) 于是由第 XI 讲 6.11 得到。b) 中各条件的\Term{equivalence}{等价}性已在第 VI$_B$ 讲 6.5.1 中证明，此处只是重述以便查考。
\end{proof}
现在可以陈述本节的主要\Term{theorem}{定理}：
\begin{SGAresult}{\Term{theorem}{定理}}{5.2}
设 $S$ 是\Term{prescheme}{预概形}，$G$ 是 $S$ 上\Term{finite-presentation}{有限呈示}的 $S$-\Term{group-prescheme}{群预概形}；以 $\mathscr L_G$（无歧义时简记为 $\mathscr L$）表示如下 $S$-\Term{functor}{函子}：对任意 $S$-\Term{prescheme}{预概形} $S'$，
\[
\mathscr L_G(S')=\left\{\begin{array}{l}
G_{S'}\text{ 的在 }S'\text{ 上光滑、纤维连通，}\\
\text{且等于自身连通正规化子的子群预概形}
\end{array}\right\}.
\]
则\Term{functor}{函子} $\mathscr L$ 可由一个 $S$-\Term{prescheme}{预概形}表示，它是一族递增的\Term{open-subprescheme}{开子预概形} $(U_i)_{i\in\mathbb N}$ 的并，各 $U_i$ 在 $S$ 上\Term{quasi-projective}{拟射影}且\Term{finite-presentation}{有限呈示}；特别地，该\Term{prescheme}{预概形}在 $S$ 上\Term{separated}{分离}。
\end{SGAresult}
\textbf{初步归约。}

对于每个整数 $r\geq0$，令 $\mathscr L_r$ 为 $\mathscr L$ 的如下\Term{subfunctor}{子函子}：对每个 $S$-\Term{prescheme}{预概形} $S'$，
\[
\mathscr L_r(S')=\left\{\begin{array}{l}
G_{S'}\text{ 的光滑、纤维连通、等于自身连通正规化子}\\
\text{且相对维数为 }r\text{ 的子群预概形}
\end{array}\right\}.
\]
由于\Term{smooth}{光滑}群的\Term{fiber}{纤维}\Term{dimension}{维数}是\Term{locally-constant}{局部常值}函数，\Term{canonical}{典范}\Term{monomorphism}{单态射} $\mathscr L_r\to\mathscr L$ 可由一个既开又闭的\Term{immersion}{浸入}表示。因此，只须对每个 $r$ 证明 $\mathscr L_r$ 可由具有以上性质的 $S$-\Term{prescheme}{预概形}表示，因为此时 $\mathscr L$ 就可由和 $S$-\Term{prescheme}{预概形} $\coprod_{r\in\mathbb N}\mathscr L_r$ 表示。

对于每个整数 $n\geq0$、每个 $S$-\Term{prescheme}{预概形} $S'$ 以及 $G_{S'}$ 的每个\Term{group-subprescheme}{子群预概形} $H$，以 $H^{(n)}$ 表示 $H$ 的\Term{unit-section}{单位截面} $S'\to H$ 的第 $n$ 个\Term{normal-invariant}{法不变量}（invariant normal，EGA IV 16.1.2）；
\SourcePage{287}
因而 $H^{(n)}$ 是一个\Term{finite-type}{有限型} $\mathcal O_{S'}$-模\Term{sheaf}{层}，对应于 $H$ 的\Term{unit-section}{单位截面}的第 $n$ 个\Term{infinitesimal-neighborhood}{无穷小邻域}。若 $H$ 在 $S'$ 上\Term{smooth}{光滑}且\Term{relative-dimension}{相对维数}为 $r$，则 $H^{(n)}$ 为\Term{locally-free}{局部自由} $\mathcal O_S$-模，其秩 $\varphi(n,r)$ 只依赖于 $n,r$。另一方面，由于 $H$ 是 $G_{S'}$ 的\Term{subprescheme}{子预概形}，有一个与\Term{base-extension}{基扩张}相容的\Term{canonical}{典范}\Term{epimorphism}{满态射}
\[G^{(n)}\otimes_{\mathcal O_S}\mathcal O_{S'}\simeq G_{S'}^{(n)}\twoheadrightarrow H^{(n)}.\]
引入\Term{projective}{射影} $S$-\Term{scheme}{概形}
\[\mathbb P_{\varphi(n,r)}=\operatorname{Grass}_{\varphi(n,r)}\bigl(G^{(n)}\bigr)\]
（EGA I，第 2 版，9.7；亦参见 Cartan 讨论班，1960/61，A.~Grothendieck 的第 14 讲）。由以上说明，映射 $H\mapsto H^{(n)}$ 定义一个\Term{canonical}{典范}\Term{morphism}{态射}
\[u_{n,r}:\mathscr L_r\longrightarrow\mathbb P_{\varphi(n,r)}.\]
群 $G$ 通过表示
\[\operatorname{int}:G\longrightarrow\operatorname{Aut}_{S\text{-gr}}(G),\qquad g\longmapsto\operatorname{int}(g),\]
自然\Term{action}{作用}于 $G^{(n)}$，因而也\Term{action}{作用}于 $\mathbb P_{\varphi(n,r)}$。此外，若 $S'$ 是\Term{quasi-compact}{拟紧}的 $S$-\Term{prescheme}{预概形}，$H\in\mathscr L_r(S')$，则我们知道（第 XI 讲 6.11），当 $n$ \Term{large-relative-to-cardinal}{足够大}时，
\[N=\operatorname{Norm}_{G_{S'}}(H)=\operatorname{Norm}_{G_{S'}}(H^{(n)}).\]
对每个整数 $n\geq0$，引入 $\mathscr L_r$ 的\Term{subfunctor}{子函子} $\mathscr L_r^n$，使对每个 $S$-\Term{prescheme}{预概形} $S'$ 有
\[\mathscr L_r^n(S')=\{H\subset G_{S'}\mid H\in\mathscr L_r(S'),\ \operatorname{Norm}_{G_{S'}}(H)=\operatorname{Norm}_{G_{S'}}(H^{(n)})\}.\]
\textbf{$\mathscr L_r^n$ 的\Term{representability}{可表性}。}

由于在 5.2 的证明结束之前，整数 $r$ 始终固定，我们在记号中省略它。因此，以 $\mathscr L,\mathscr L^n,\mathbb P_n,u_n$ 分别代替 $\mathscr L_r,\mathscr L_r^n,\mathbb P_{\varphi(n,r)},u_{n,r}$。

令 $v_n$ 为 $u_n$ 在 $\mathscr L$ 的\Term{subfunctor}{子函子} $\mathscr L^n$ 上的\Term{restriction}{限制}。根据 $\mathscr L^n$ 的定义及 5.1 b) ii)，$v_n$ 是\Term{monomorphism}{单态射}。事实上，有以下\Term{lemma}{引理}：
\begin{SGAresult}{\Term{lemma}{引理}}{5.3}
\Term{morphism}{态射} $v_n$ 是\Term{finite-presentation}{有限呈示}\Term{immersion}{浸入}。特别地，$\mathscr L^n$ 可由一个在 $S$ 上\Term{quasi-projective}{拟射影}且\Term{finite-presentation}{有限呈示}的 $S$-\Term{prescheme}{预概形}表示。
\end{SGAresult}
用 $\mathbb P_n$ 替换 $S$，通过通常的技巧，归结为证明以下断言：设 $Q\in\mathbb P_n(S)$，考察由 $\operatorname{Sch}/S$ 的\Term{final-object}{终对象} $S$ 表示的\Term{functor}{函子} $h_S$ 的\Term{subfunctor}{子函子} $\mathscr F$，其定义为：对于每个 $S$-\Term{prescheme}{预概形} $S'$，
\SourcePage{288}
\[\mathscr F(S')=\begin{cases}
h_S(S')\text{（单元素集合）},&\text{若存在 }H\in\mathscr L^n(S')\text{ 使 }H^{(n)}=Q_{S'},\\
\varnothing,&\text{否则}.
\end{cases}\]
则\Term{canonical}{典范}\Term{monomorphism}{单态射} $\mathscr F\to S$ 是\Term{finite-presentation}{有限呈示}\Term{immersion}{浸入}。

先把\Term{functor}{函子} $\mathscr F$ 的定义改写。注意，$Q$ 在 $G$ 中的\Term{normalizer}{正规化子}可由一个在 $S$ 上\Term{finite-presentation}{有限呈示}的\Term{group-subprescheme}{子群预概形} $N$ 表示，即 $\mathbb P_n(S)$ 的点 $Q$ 在\Term{morphism}{态射}
\[G\longrightarrow\mathbb P_n,\qquad g\longmapsto g\cdot Q\]
下的\Term{inverse-image}{逆像}。我们断言，\Term{functor}{函子} $\mathscr F$ 等于 $h_S$ 的如下\Term{subfunctor}{子函子}：若满足

(a) $N_{S'}$ 沿\Term{unit-section}{单位截面}\Term{smooth}{光滑}且\Term{relative-dimension}{相对维数}为 $r$；

(b) $(N_{S'})^{(n)}$（此时\Term{canonical}{典范}地是 $\mathbb P_n(S')$ 的一个元素）等于 $Q_{S'}$，

则 $\mathscr F(S')=h_S(S')$；否则，$\mathscr F(S')=\varnothing$。

事实上，暂以 $\mathscr F_1$ 表示第一种定义的\Term{functor}{函子}，$\mathscr F_2$ 表示第二种定义的\Term{functor}{函子}。于是：
\[\text{i)}\qquad\mathscr F_1(S')=h_S(S')\ \Longrightarrow\ \mathscr F_2(S')=h_S(S').\]
事实上，令 $H\in\mathscr F^n(S')$ 为 $G_{S'}$ 中满足 $H^{(n)}=Q_{S'}$ 的\Term{subgroup}{子群}。于是
\[\operatorname{Norm}_{G_{S'}}(H)=\operatorname{Norm}_{G_{S'}}(H^{(n)})=\operatorname{Norm}_{G_{S'}}(Q_{S'})=N_{S'}.\]
所以，根据 5.1 b) ii)，$N_{S'}$ 沿\Term{unit-section}{单位截面}\Term{smooth}{光滑}，其\Term{identity-component}{单位连通分支}为 $H$。从而 $N_{S'}$ \Term{relative-dimension}{相对维数}为 $r$，且因 $H$ 是 $N_{S'}$ 中的\Term{open-subset}{开集}（5.1 b) i)），有 $(N_{S'})^{(n)}=H^{(n)}=Q_{S'}$。总之，$\mathscr F_2(S')=h_S(S')$。
\[\text{ii)}\qquad\mathscr F_2(S')=h_S(S')\ \Longrightarrow\ \mathscr F_1(S')=h_S(S').\]
根据假设，$N_{S'}$ 沿\Term{unit-section}{单位截面}\Term{smooth}{光滑}，\Term{relative-dimension}{相对维数}为 $r$；其\Term{identity-component}{单位连通分支}因而可由一个\Term{group-subprescheme}{子群预概形} $H$ 表示（第 VI$_B$ 讲 3.10），$H$ 在 $S'$ 上\Term{smooth}{光滑}，\Term{fiber}{纤维}\Term{connected}{连通}且\Term{dimension}{维数}为 $r$。由于 $H$ 在 $N_{S'}$ 中\Term{normal-subgroup}{正规}且为\Term{open-subset}{开集}，有下列包含关系：
\begin{align*}
N_{S'}&\subset\operatorname{Norm}_{G_{S'}}(H)\subset\operatorname{Norm}_{G_{S'}}(H^{(n)})\\
&=\operatorname{Norm}_{G_{S'}}\bigl(N_{S'}^{(n)}\bigr)=\operatorname{Norm}_{G_{S'}}(Q_{S'})=N_{S'}.
\end{align*}
所以这些包含关系都是等号。第一处包含关系于是说明 $H$ 等于自身的\Term{connected-normalizer}{连通正规化子}，第二处说明 $H$ 是 $\mathscr L^n(S')$ 的一个元素。也就是说，$\mathscr F_1(S')=h_S(S')$。

i) 和 ii) 两个蕴含关系给出 $\mathscr F_1=\mathscr F_2$。保留\Term{functor}{函子} $\mathscr F$ 的第二种定义，先用一个\Term{finite-presentation}{有限呈示}\Term{immersion}{浸入}来“表示条件 a)”。为此，只须应用：
\SourcePage{289}
\begin{SGAresult}{\Term{lemma}{引理}}{5.4}
设 $S$ 是\Term{prescheme}{预概形}，$X$ 是 $S$ 上\Term{locally-finite-presentation}{局部有限呈示}的 $S$-\Term{prescheme}{预概形}，$\sigma:S\to X$ 是 $X$ 的一个\Term{section}{截面}，$r\geq0$ 是整数，$L:(\operatorname{Sch}/S)^{\circ}\to\operatorname{Ens}$ 是 $S$ 的如下\Term{subfunctor}{子函子}：若 $X_{S'}$ 沿\Term{section}{截面} $\sigma_{S'}$ \Term{smooth}{光滑}，且在 $\sigma_{S'}(S')$ 的各点处\Term{relative-dimension}{相对维数}为 $r$，则 $L(S')=h_S(S')$；否则 $L(S')=\varnothing$。则：

a) \Term{monomorphism}{单态射} $L\to S$ 是\Term{finite-presentation}{有限呈示}\Term{immersion}{浸入}。

b) 令 $J$ 为\Term{immersion}{浸入} $S\to X$ 的\Term{conormal-sheaf}{余法层}（EGA IV 16.1.2）；若对于 $S$ 的每一点 $s$，$J\otimes_{\mathcal O_{S,s}}\kappa(s)$ 的秩至多为 $r$，则\Term{immersion}{浸入} $L\to S$ 是\Term{closed-immersion}{闭浸入}。
\end{SGAresult}
\begin{proof}
\Term{functor}{函子} $L$ 在 $S$ 上具有\Term{local-character}{局部性}，这使我们归结到 $S$ \Term{affine}{仿射}的情形。用 $\sigma(S)$ 的一个\Term{neighborhood}{邻域}替换 $X$，可以假定 $X$ 在 $S$ 上\Term{finite-presentation}{有限呈示}，再由 EGA IV 8.9 假定 $S$ 为 noether（在 b) 的情形中，注意\Term{conormal-sheaf}{余法层}的构造与\Term{base-extension}{基扩张}交换（EGA IV 16.6.4），而 $J$ 的\Term{fiber}{纤维}的秩是 $S$ 上的\Term{constructible}{可构造}函数）。对于每个 $S$-\Term{prescheme}{预概形} $S'$，以 $J'=J_{S'}$ 表示\Term{section}{截面} $\sigma_{S'}$ 的\Term{conormal-sheaf}{余法层}；令 $\operatorname S^{\bullet}(J')$ 为 $J'$ 在 $\mathcal O_{S'}$ 上的\Term{symmetric-algebra}{对称代数}，它与 $\operatorname S^{\bullet}(J)_{S'}$ \Term{canonical-isomorphism}{典范同构}；令 $\operatorname{Gr}_{\bullet}(\sigma_{S'})$ 为与 $\sigma_{S'}$ 相关联的\Term{graded}{分次} $\mathcal O_{S'}$-\Term{algebra}{代数}\Term{sheaf}{层}（EGA IV 16）；对每个整数 $n\geq0$，令 $\sigma_{n,S'}:\operatorname S^n(J')\to\operatorname{Gr}_n(\sigma_{S'})$ 为\Term{canonical}{典范}\Term{epimorphism}{满态射}。

由 EGA IV 17.12.3 及 EGA 0$_{\mathrm{IV}}$ 19.5.4，$X_{S'}$ 沿\Term{section}{截面} $\sigma_{S'}$ \Term{smooth}{光滑}，且在 $\sigma_{S'}(S')$ 的各点处\Term{relative-dimension}{相对维数}为 $r$，当且仅当：

i) $J'$ 是秩为 $r$ 的\Term{locally-free}{局部自由} $\mathcal O_{S'}$-\Term{sheaf}{层}。

ii) 对每个整数 $n\geq1$，$\sigma_{n,S'}$ 都是\Term{isomorphism}{同构}。

由 TDTE IV，\Term{lemma}{引理} 3.6，“使 $J$ \Term{locally-free}{局部自由}且秩为 $r$ 的\Term{functor}{函子}”可由一个\Term{subprescheme}{子预概形} $S_1$ 表示，并且在 b) 的情形中 $S_1$ 在 $S$ 中闭。用 $S_1$ 替换 $S$，便归结到 $J$ \Term{locally-free}{局部自由}且秩为 $r$ 的情形。对整数 $n$ 进行归纳。假设已用 $S$ 的一个\Term{closed-subprescheme}{闭子预概形} $S_{n-1}$ 表示了“对每个整数 $q\leq n-1$，使\Term{morphism}{态射} $\sigma_{q,S'}$ 为\Term{injectivity}{单射}的\Term{functor}{函子}”；下面证明“使 $\sigma_{n,S'}$ 为\Term{injectivity}{单射}的\Term{subfunctor}{子函子}”\footnote{编者注：这里以“\Term{subfunctor}{子函子}”替换了“\Term{functor}{函子}”。是否应写成“对 $q\leq n$ 使 $\sigma_{q,S'}$ 为\Term{injectivity}{单射}”？}可由 $S_{n-1}$ 的一个\Term{closed-subprescheme}{闭子预概形} $S_n$ 表示。用 $S_{n-1}$ 替换 $S$，可以假定对 $q\leq n-1$，$\sigma_{q,S}$ 为\Term{bijection}{双射}。但此时，关于\Term{section}{截面} $\sigma_s$ 的第 $(n-1)$ 个\Term{normal-invariant}{法不变量} $X^{(n-1)}$ 有一个\Term{composition-series}{合成列} $X^{(0)},\ldots,X^{(n-1)}$，其相继商 $\operatorname{Gr}_0(\sigma_s),\ldots,\operatorname{Gr}_{n-1}(\sigma_s)$ 都是\Term{locally-free}{局部自由} $\mathcal O_{S'}$-模，故\Term{flat}{平坦}，从而 $X^{(n-1)}$ 在 $\mathcal O_S$ 上\Term{flat}{平坦}。由于 $X^{(i)}$ 的构造与\Term{base-extension}{基扩张}交换（EGA IV 16），对每个 $S$-\Term{prescheme}{预概形} $S'$，有
\[\operatorname{Gr}_n(\sigma_{S'})=\operatorname{Ker}\bigl(X_{S'}^{(n)}\to X_{S'}^{(n-1)}\bigr)=(\operatorname{Gr}_n(\sigma))_{S'},\qquad \sigma_{n,S'}=(\sigma_{n,S})_{S'}.\]
\SourcePage{290}
所以，我们关注的\Term{functor}{函子}就是“使\Term{morphism}{态射} $\sigma_{n,S}$ 为\Term{injectivity}{单射}的\Term{functor}{函子}”。它在 $S$ 上具有\Term{local-character}{局部性}，因此可以假定 $S$ \Term{affine}{仿射}，且 $\operatorname S^n(J)$ 是自由 $\mathcal O_S$-模。此时显然，这个\Term{functor}{函子}可由 $S$ 的一个\Term{closed-subscheme}{闭子概形}表示，其\Term{ideal-of-definition}{定义理想}由 $\operatorname{Ker}\sigma_{n,S}$ 关于 $\operatorname S^n(J)$ 的一组基的坐标生成。

由于 $S$ 为 noether，$S_1$ 的\Term{closed-subprescheme}{闭子预概形}递减序列 $\{S_n\}$ \Term{stable}{稳定}；其\Term{stable}{稳定}值表示\Term{functor}{函子} $L$，这就完成了 5.4 的证明。
\end{proof}
回到 $\mathscr L^n$ 的\Term{representability}{可表性}问题。用适当的\Term{subprescheme}{子预概形} $S_1$ 替换 $S$，可以假定 $N$ 沿\Term{unit-section}{单位截面}\Term{smooth}{光滑}，且\Term{relative-dimension}{相对维数}为 $r$。此时\Term{functor}{函子} $\mathscr L^n$ 就是 $\mathbb P_n$ 在 $S$ 上的两个\Term{section}{截面} $h,g$ 的“相等\Term{functor}{函子}”，它们分别对应于\Term{sheaf}{层} $Q$ 及 $N^{(n)}$（即上面\Term{functor}{函子} $\mathscr F_2$ 定义中的条件 b)）。因此，它可由 $S$ 的\Term{closed-subprescheme}{闭子预概形}表示，即 $\mathbb P_n\times_S\mathbb P_n$ 的对角在\Term{finite-presentation}{有限呈示}\Term{morphism}{态射} $h\times_Sg$ 下的\Term{inverse-image}{逆像}。5.3 的证明至此完成。

\textbf{\Term{transition-morphism}{过渡态射} $\mathscr L^n\to\mathscr L^m$ 的研究。}

若 $G_{S'}$ 的一个\Term{subgroup}{子群} $H$ 属于 $\mathscr L^n(S')$，则对于每个 $m\geq n$，它也属于 $\mathscr L^m(S')$；由此得到自然\Term{monomorphism}{单态射}
\[u_n^m:\mathscr L^n\longrightarrow\mathscr L^m\qquad(m\geq n).\]
\begin{SGAresult}{\Term{lemma}{引理}}{5.5}
\Term{morphism}{态射} $u_n^m$ 是\Term{open-immersion}{开浸入}。
\end{SGAresult}
改变 $S$ 后，归结为以下问题：设 $H\in\mathscr L^m(S)$，
\[N=\operatorname{Norm}_G(H)=\operatorname{Norm}_G(H^{(m)}),\qquad N'=\operatorname{Norm}_G(H^{(n)}),\]
令 $\mathscr D:(\operatorname{Sch}/S)^{\circ}\to\operatorname{Ens}$ 为 $N,N'$ 的相等\Term{functor}{函子}，定义为：若 $N_{S'}=N'_{S'}$，则 $\mathscr D(S')=h_S(S')$，否则 $\mathscr D(S')=\varnothing$。须证明 $\mathscr D\to S$ 是\Term{open-immersion}{开浸入}。我们断言，$\mathscr D$ 也就是 $S$ 的“使\Term{immersion}{浸入} $H\to N'$ 成为\Term{open-immersion}{开浸入}”的\Term{subfunctor}{子函子}。事实上，若 $N_{S'}=N'_{S'}$，则 $H_{S'}\to N'_{S'}$ 确为\Term{open-immersion}{开浸入}，因为 $H_{S'}\to N_{S'}$ 如此（命题 5.1）。反之，若 $H_{S'}\to N'_{S'}$ 是\Term{open-immersion}{开浸入}，由于 $H$ 的\Term{fiber}{纤维}\Term{connected}{连通}，$H_{S'}$ 就是 $N'_{S'}$ 的\Term{identity-component}{单位连通分支}（第 VI$_B$ 讲 3.10），从而在 $N'_{S'}$ 中\Term{normal-subgroup}{正规}，故 $N'_{S'}\subset N_{S'}$。而 $N'$ 总包含 $N$，所以确有 $N_{S'}=N'_{S'}$。\Term{group-prescheme}{群预概形} $H,N'$ 均在 $S$ 上\Term{finite-presentation}{有限呈示}，且 $H$ 在 $S$ 上\Term{flat}{平坦}；$\mathscr D\to S$ 是\Term{open-immersion}{开浸入}这一事实便由第 VI$_B$ 讲 2.6 得到。

\textbf{完成 5.2 的证明。}

由于各\Term{functor}{函子} $\mathscr L^n$ \Term{representability}{可表}，且\Term{transition-morphism}{过渡态射} $u_n^m$ 相容并可由\Term{open-immersion}{开浸入}表示，存在一个 $S$-\Term{prescheme}{预概形} $X$，为一列递增\Term{open-subset}{开集} $(X_i)_{i\in\mathbb N}$ 的并，其中 $X_i$ 表示\Term{functor}{函子} $\mathscr L^i$，且把 $\mathscr L^i$ 与 $X_i$ 等同后，\Term{inclusion-morphism}{包含态射} $X_i\to X_j$（$j\geq i$）就是 $u_i^j$。为\Term{pushout}{推出} $X$ 表示\Term{functor}{函子} $\mathscr L$，只须注意：
\SourcePage{291}
在赋予 Zariski \Term{topology}{拓扑}的 $\operatorname{Sch}/S$ 上的\Term{sheaf}{层}\Term{category}{范畴}中（第 IV 讲 6.1），有
\[X=\varprojlim X_i\qquad\text{及}\qquad\mathscr L=\varprojlim\mathscr L^i.\]
\begin{SGAresult}{注记}{5.6}
采用以上记号，进一步假设 $S$ 的所有\Term{residue-characteristic}{剩余特征}均为零，则对每个整数 $r\geq0$，\Term{functor}{函子} $\mathscr L_r$ 等于 $\mathscr L_r^1$，从而可由一个在 $S$ 上\Term{finite-presentation}{有限呈示}且\Term{quasi-projective}{拟射影}的 $S$-\Term{prescheme}{预概形}表示。事实上，只须证明：若 $H\in\mathscr L_r(S')$，则\Term{canonical}{典范}\Term{immersion}{浸入}
\[H\longrightarrow N=\operatorname{Norm}_{G_{S'}}(H^{(1)})\]
是\Term{open-immersion}{开浸入}（因为这就\Term{pushout}{推出} $N=\operatorname{Norm}_{G_{S'}}(H)$，故 $H\in\mathscr L_r^1(S')$）。由于 $H$ 在 $S$ 上\Term{flat}{平坦}，$H,N$ 在 $S$ 上\Term{finite-presentation}{有限呈示}，只须证明（第 VI$_B$ 讲 2.6），对于 $S$ 的每一点 $s$，$H_s\to N_s$ 是\Term{open-immersion}{开浸入}。而由 Cartier \Term{theorem}{定理}（第 VI$_B$ 讲 1.6）容易\Term{pushout}{推出}\footnote{编者注：例如可以应用 M.~Demazure 与 P.~Gabriel，\emph{Groupes algébriques I}，Masson \& North Holland（1970），命题 II.6.1。}，若 $G$ 是\Term{characteristic-zero}{特征零}的域上的\Term{algebraic-group}{代数群}，$H$ 是\Term{connected}{连通}\Term{algebraic-subgroup}{代数子群}，则
\[\operatorname{Norm}_G(H)=\operatorname{Norm}_G(\operatorname{Lie}H)=\operatorname{Norm}_G(H^{(1)}).\]
反之，若 $S$ 具有非零\Term{residue-characteristic}{剩余特征}，$\mathscr L_r$ 的\Term{subfunctor}{子函子} $\mathscr L_r^n$ 可能构成严格递增的序列（即使 $S$ \Term{quasi-compact}{拟紧}），此时 $\mathscr L_r$ 不能由一个在 $S$ 上\Term{quasi-compact}{拟紧}的 $S$-\Term{prescheme}{预概形}表示。例如，取\Term{characteristic}{特征} $p>0$ 的域 $k$ 上的\Term{algebraic-group}{代数群} $G$，它是\Term{torus}{环面} $T=\mathbb G_m\times\mathbb G_m$ 与\Term{unipotent-group}{幺幂群} $U=\mathbb G_a\times\mathbb G_a$ 的\Term{semidirect-product}{半直积}，其中 $T$ 在 $U$ 上的\Term{action}{作用}定义为 $(t,t',u,u')\mapsto(tu,t'u')$。对每个整数 $n>0$，考察 $U$ 中由方程 $u'=u^{p^n}$ 定义的\Term{smooth}{光滑}\Term{connected}{连通}\Term{subgroup}{子群} $U_n$，以及 $T$ 中由方程 $t'=t^{p^n}$ 定义的\Term{subtorus}{子环面} $T_n$。立即可见，$T_n$ \Term{action}{作用}于 $U_n$，且 $G$ 中等于 $T_n\cdot U_n$ 的\Term{subgroup}{子群} $G_n$ \Term{smooth}{光滑}、\Term{connected}{连通}，并等于其在 $G$ 中的\Term{normalizer}{正规化子}。然而，所有 $n\geq m$ 对应的群 $G_n$ 互不相同，却具有相同的 $p^m$ 阶\Term{infinitesimal-neighborhood}{无穷小邻域}。
\end{SGAresult}
\begin{SGAresult}{注记}{5.7}
$\mathscr L$ 上有一个\Term{canonical}{典范}\Term{invertible-sheaf}{可逆层} $L$，它在 $\mathscr L$ 的任一对 $S$ \Term{quasi-compact}{拟紧}的\Term{open-subset}{开集} $U$ 上的\Term{restriction}{限制}都为 $S$-\Term{ample}{丰沛}。事实上，考察 $G_{\mathscr L}$ 的\Term{group-subprescheme}{子群预概形} $H$，它在 $\mathscr L$ 上\Term{smooth}{光滑}，\Term{fiber}{纤维}\Term{connected}{连通}，等于自身的\Term{connected-normalizer}{连通正规化子}，并对这些性质具有\Term{versal}{万有}性。我们断言，可以取 $L=(\det(\operatorname{Lie}H))^{-1}$（回忆：若 $F$ 是\Term{prescheme}{预概形} $S$ 上的一个\Term{finite-rank}{有限秩}\Term{locally-free}{局部自由} $\mathcal O_S$-模\Term{sheaf}{层}，则 $\det(F)$ 表示这样一个\Term{invertible}{可逆} $\mathcal O_S$-模：在 $F$ 秩为 $r$ 的既开又闭\Term{subprescheme}{子预概形} $S_r$（$r\geq0$）上的\Term{restriction}{限制}为 $\bigwedge^r(F)$）。沿用 5.2 证明中的记号。为证明关于 $L$ 的断言，只须限于\Term{functor}{函子} $\mathscr L_r^n$，并证明 $L|_{\mathscr L_r^n}$ 为 $S$-\Term{ample}{丰沛}。考察\Term{canonical}{典范}\Term{immersion}{浸入} $v_n:\mathscr L_r^n\to\mathbb P_n$，令 $Q$ 为 $\mathbb P_n$ 上对\Term{grassmannian}{Grassmann 概形} $\mathbb P_n$ 具有\Term{versal}{万有}性的\Term{locally-free}{局部自由}\Term{sheaf}{层}。按构造有 $v_n^*(Q)=H^{(n)}$（此时 $H$ 表示 $G_{\mathscr L_r^n}$ 中对\Term{functor}{函子} $\mathscr L_r^n$ 具有\Term{versal}{万有}性的\Term{group-subprescheme}{子群预概形}）。而 $\det(Q)$ 是 $\mathbb P_n$ 上的\Term{canonical}{典范}\Term{ample}{丰沛}\Term{sheaf}{层}
\SourcePage{292}
（EGA I，第 2 版，9.7），因此 $\det H^{(n)}$ 对 $\mathscr L_r^n$ 相对\Term{ample}{丰沛}（EGA II 4.6.13 i) bis）。仍以 $J$ 表示\Term{conormal-sheaf}{余法层}，它等于 $H^{(1)}$，以 $\operatorname S^q(J)$ 表示 $J$ 在 $\mathcal O_{\mathscr L_r^n}$ 上的\Term{symmetric-algebra}{对称代数}的 $q$ 次\Term{homogeneous-component}{齐次部分}。由于 $H$ 在 $\mathscr L_r^n$ 上\Term{smooth}{光滑}，立即得到\Term{canonical-isomorphism}{典范同构}
\[\det H^{(n)}\simeq\prod_{1\leq q\leq n}\det\operatorname S^q(J).\]
另一方面，可以证明，对每个秩为 $r$ 的\Term{locally-free}{局部自由}\Term{sheaf}{层} $J$ 及每个整数 $q>0$，都有\Term{canonical-isomorphism}{典范同构}
\[\det\operatorname S^q(J)\simeq(\det J)^{\otimes s},\]
其中整数 $s>0$ 只依赖于 $r,q$。最后，对某个适当的整数 $s>0$，得到 $\det H^{(n)}\simeq(\det J)^{\otimes s}$，所以（EGA II 4.5.6），$\det J=(\det\operatorname{Lie}H)^{-1}$ 确为 $S$-\Term{ample}{丰沛}。
\end{SGAresult}
\begin{SGAresult}{注记}{5.8}
设 $S$ 是\Term{prescheme}{预概形}，$G,H$ 是两个在 $S$ 上\Term{finite-presentation}{有限呈示}的 $S$-\Term{group-prescheme}{群预概形}，$i:H\to G$ 是一个作为\Term{monomorphism}{单态射}的 $S$-\Term{group-homomorphism}{群同态}。若 $H$ 在 $S$ 上\Term{smooth}{光滑}且\Term{fiber}{纤维}\Term{connected}{连通}，我们知道（第 XI 讲 6.11），$N=\operatorname{Norm}_G(H)$ 可由 $G$ 的一个在 $S$ 上\Term{finite-presentation}{有限呈示}的\Term{closed-group-subprescheme}{闭子群预概形}表示。再假设 $N$ 沿\Term{unit-section}{单位截面}\Term{smooth}{光滑}，且在 $S$ 上与 $H$ \Term{relative-dimension}{相对维数}相同。此时，$N$ 的\Term{identity-component}{单位连通分支} $N^0$ 可由 $N$ 的一个在 $S$ 上\Term{smooth}{光滑}的\Term{open-group-subprescheme}{开子群预概形}表示（第 VI$_B$ 讲 3.10）。\Term{monomorphism}{单态射} $i$ 显然经由 $N^0$ \Term{resolution}{分解}。事实上，$H=N^0$。因为对 $S$ 的每一点 $s$，都有 $H_s=(N^0)_s$：这两个\Term{algebraic-group}{代数群}均\Term{connected}{连通}、\Term{smooth}{光滑}且\Term{dimension}{维数}相同。由于 $H$ 在 $S$ 上\Term{flat}{平坦}，由此\Term{pushout}{推出} $H\to N^0$ 是\Term{isomorphism}{同构}（EGA IV 17.9.5）。最终，$H$ 是 $G$ 的一个\Term{group-subprescheme}{子群预概形}。所以，本节引入的\Term{functor}{函子} $\mathscr L$ 就是由 $G$ 中那些在 $S$ 上\Term{smooth}{光滑}、\Term{fiber}{纤维}\Term{connected}{连通}并等于自身\Term{connected-normalizer}{连通正规化子}的\Term{subgroup}{子群} $H$ 所成的\Term{functor}{函子}。
\end{SGAresult}
\SGAsection{6}{Cartan 子群函子与抛物子群函子}
当 $G$ 是\Term{algebraically-closed-field}{代数闭域} $k$ 上的\Term{smooth}{光滑}\Term{connected}{连通}\Term{algebraic-group}{代数群}时，我们已经定义了 $G$ 的\Term{subtorus}{子环面}、\Term{maximal-torus}{极大子环面}、\Term{cartan-subgroup}{Cartan 子群}（第 XII 讲 1）、\Term{borel-subgroup}{Borel 子群}（第 XIV 讲 4.1）、\Term{parabolic-subgroup}{抛物子群}（第 XIV 讲 4.8 bis）。现在按如下方式，将这些概念推广到任意\Term{base}{基底}上的\Term{group-prescheme}{群预概形}：
\begin{SGAresult}{定义}{6.1}
设 $S$ 是\Term{prescheme}{预概形}，$G$ 是 $S$ 上\Term{finite-presentation}{有限呈示}的 $S$-\Term{group-prescheme}{群预概形}，$H$ 是 $S$-\Term{group-prescheme}{群预概形}，$i:H\to G$ 是一个使 $H$ 成为 $G$ 的\Term{subgroup}{子群}的 $S$-\Term{monomorphism}{单态射}。若满足下列条件，就称 $H$ 是 $G$ 的\Term{subtorus}{子环面}（相应地，\Term{maximal-torus}{极大子环面}、\Term{cartan-subgroup}{Cartan 子群}、\Term{borel-subgroup}{Borel 子群}、\Term{parabolic-subgroup}{抛物子群}）：

i) $H$ 在 $S$ 上\Term{smooth}{光滑}。
\SourcePage{293}
ii) 对 $S$ 上的每个\Term{geometric-point}{几何点} $\bar s$，$H_{\bar s}$ 是 $(G_{\bar s})_{\mathrm{red}}^0$ 的\Term{subtorus}{子环面}（相应地，\Term{maximal-torus}{极大子环面}、\Term{cartan-subgroup}{Cartan 子群}、\Term{borel-subgroup}{Borel 子群}、\Term{parabolic-subgroup}{抛物子群}）。
\end{SGAresult}
\begin{SGAresult}{注记}{6.1 bis}
a) 若 $S$-群 $H$ 是 $G$ 的\Term{subtorus}{子环面}（相应地，其余各情形），则其\Term{fiber}{纤维}\Term{connected}{连通}，从而 $H$ 在 $S$ 上\Term{finite-presentation}{有限呈示}（第 VI$_B$ 讲 5.3.3）。

b) 若 $H$ 是 $G$ 的\Term{subtorus}{子环面}，则 $H$ 是第 IX 讲意义下的\Term{torus}{环面}，这立即由第 X 讲 8.1 得到。此外，\Term{monomorphism}{单态射} $i:H\to G$ 是\Term{immersion}{浸入}（参见下文 8.3）。

c) 若 $G$ 在 $S$ 上\Term{smooth}{光滑}且\Term{fiber}{纤维}\Term{connected}{连通}，$H$ 是 $G$ 的\Term{cartan-subgroup}{Cartan 子群}（相应地，\Term{borel-subgroup}{Borel 子群}、\Term{parabolic-subgroup}{抛物子群}），则\Term{monomorphism}{单态射} $i:H\to G$ 是\Term{immersion}{浸入}，因此这里的定义与第 XII、XIV 讲引入的定义一致。事实上，此时 $H$ 等于自身的\Term{connected-normalizer}{连通正规化子}（第 XII 讲 6.6 c)，第 XIV 讲 4.8 及 4.8 bis），只须应用 5.8。
\end{SGAresult}
\begin{SGAresult}{定义}{6.1 ter}
设 $S$ 是\Term{prescheme}{预概形}，$G$ 是\Term{locally-finite-type}{局部有限型} $S$-\Term{group-prescheme}{群预概形}，$s$ 是 $S$ 的一点。称 $(G_{\bar s})_{\mathrm{red}}^0$ 的\Term{cartan-subgroup}{Cartan 子群}的\Term{dimension}{维数}为 $G$ 在点 $s$ 处的\Term{nilpotent-rank}{幂零秩}，记为 $\rho_{\mathrm n}(s)$。类似地，定义\Term{reductive-rank}{约化秩} $\rho_{\mathrm r}(s)$、\Term{unipotent-rank}{幺幂秩} $\rho_{\mathrm u}(s)$、\Term{abelian-rank}{阿贝尔秩} $\rho_{\mathrm{ab}}(s)$（参见第 X 讲 8.7）。
\end{SGAresult}
现设 $G$ 是\Term{algebraically-closed-field}{代数闭域} $k$ 上的\Term{smooth}{光滑}\Term{connected}{连通}\Term{algebraic-group}{代数群}。回忆 $G$ 的\Term{radical}{根基}，记为 $\operatorname{rad}(G)$，是 $G$ 的最大的\Term{normal-subgroup}{正规}、\Term{smooth}{光滑}、\Term{connected}{连通}且\Term{solvable}{可解}的\Term{algebraic-subgroup}{代数子群}；此时 $G/\operatorname{rad}(G)$ 是\Term{semisimple}{半单}的（利用第 XII 讲 6.1 归结到 $G$ \Term{affine}{仿射}的情形）。若 $G$ 还\Term{affine}{仿射}，则定义 $G$ 的\Term{unipotent-radical}{幺幂根基} $\operatorname{rad}^{\mathrm u}(G)$ 为 $G$ 的最大的\Term{normal-subgroup}{正规}、\Term{smooth}{光滑}、\Term{connected}{连通}且\Term{unipotent}{幺幂}的\Term{algebraic-subgroup}{代数子群}；此时 $G/\operatorname{rad}^{\mathrm u}(G)$ 是\Term{reductive}{约化}的。
\begin{SGAresult}{命题}{6.2}
设 $S$ 是\Term{prescheme}{预概形}，$G,H$ 是两个在 $S$ 上\Term{finite-presentation}{有限呈示}的 $S$-\Term{group-prescheme}{群预概形}，$i:H\to G$ 是使 $H$ 成为 $G$ 的\Term{subgroup}{子群}的 $S$-群\Term{monomorphism}{单态射}，$\mathrm P(s)$ 是关于 $S$ 的点 $s$ 的下列性质之一：

i) $(G_{\bar s})_{\mathrm{red}}^0$ 是\Term{abelian-variety}{阿贝尔簇}（相应地，\Term{affine}{仿射}、是\Term{torus}{环面}、\Term{unipotent}{幺幂}）。

ii) $(H_{\bar s})_{\mathrm{red}}^0$ 是 $G_{\bar s}$ 的\Term{maximal-torus}{极大环面}。

iii) $(H_{\bar s})_{\mathrm{red}}^0$ 是 $G_{\bar s}$ 的某个\Term{torus}{环面}在 $(G_{\bar s})_{\mathrm{red}}^0$ 中的\Term{centralizer}{中心化子}（相应地，是 $G_{\bar s}$ 的\Term{cartan-subgroup}{Cartan 子群}）。

iv) $(H_{\bar s})_{\mathrm{red}}^0$ 是 $G_{\bar s}$ 的\Term{borel-subgroup}{Borel 子群}（相应地，\Term{parabolic-subgroup}{抛物子群}）。

v) $(H_{\bar s})_{\mathrm{red}}^0$ 是 $G_{\bar s}$ 的\Term{radical}{根基}（相应地，$(G_{\bar s})_{\mathrm{red}}^0$ 是\Term{semisimple}{半单}的）。

vi) $G_{\bar s}$ \Term{affine}{仿射}且 $(H_{\bar s})_{\mathrm{red}}^0$ 是 $G_{\bar s}$ 的\Term{unipotent-radical}{幺幂根基}（相应地，$(G_{\bar s})_{\mathrm{red}}^0$ 是\Term{reductive}{约化}的）。

则 $S$ 中满足 $\mathrm P(s)$ 的点 $s$ 所成的\Term{set}{集合} $E$ 是\Term{locally-constructible}{局部可构造}的（EGA 0$_{\mathrm{III}}$ 9.1.11）。
\end{SGAresult}
\SourcePage{294}
\begin{SGAresult}{注记}{6.2.1}
这个命题补充了第 VI$_B$ 讲第 10 节。此外，利用\Term{semicontinuity-theorem}{半连续性定理}（参见第 X 讲 8.7），还可以进一步说明 $E$ 的结构；稍后将看到一个例子。
\end{SGAresult}
\textbf{6.2 的证明。}

注意，如果 $S$ 是域的谱，则 $E$ 在该域的\Term{group-extension}{扩张}下\Term{invariant}{不变}。通过标准归约（EGA IV 9），可以归结到 $S$ 是\Term{generic-point}{泛点}为 $\eta$ 的 noether \Term{integral}{整}\Term{prescheme}{预概形}的情形。须证明 $E$ 或 $\operatorname{ens}(S)\setminus E$ 包含 $\eta$ 的一个\Term{neighborhood}{邻域}（EGA IV 9.2.1）。可以假定 $S$ \Term{affine}{仿射}，其环为 $A$，\Term{field-of-fractions}{分式域}为 $K$。若 $L$ 是 $K$ 的\Term{finite-extension}{有限扩张}，显然存在 $L$ 的一个在 $A$ 上\Term{finite}{有限}且以 $L$ 为\Term{field-of-fractions}{分式域}的子 $A$-\Term{algebra}{代数} $B$。置 $S'=\operatorname{Spec}B$，则\Term{canonical}{典范}\Term{morphism}{态射} $S'\to S$ \Term{dominant}{支配}且\Term{finite-presentation}{有限呈示}，所以 $S'$ 中的一个\Term{nonempty}{非空}\Term{open-subset}{开集}的像包含 $S$ 的一个\Term{nonempty}{非空}\Term{open-subset}{开集}（EGA IV 1.8.4）。对于当前的问题，可以用 $S'$ 替换 $S$，从而用一个\Term{finite-extension}{有限扩张} $L$ 替换 $K$。因此，可以选择 $L$，使 $(G_L)_{\mathrm{red}}$ 及 $(H_L)_{\mathrm{red}}$ 在 $L$ 上\Term{smooth}{光滑}（EGA IV 4.6.6）。缩小 $S'$ 后，可以假定 $G_{\mathrm{red}},H_{\mathrm{red}}$ 是 $S$ 上的\Term{smooth}{光滑}\Term{group-prescheme}{群预概形}（第 VI$_B$ 讲第 10 节及 EGA IV 17）。鉴于要证明的性质，可以用 $G,H$ 的\Term{reduced}{约化}\Term{identity-component}{单位连通分支}替换它们（第 VI$_B$ 讲 10.9），从而假定 $G,H$ 在 $S$ 上\Term{smooth}{光滑}，且\Term{fiber}{纤维}\Term{connected}{连通}。最后，可以假定 $H$ 是 $G$ 的\Term{closed-group-subprescheme}{闭子群预概形}（第 VI$_B$ 讲 10.4）。

\emph{i) 的证明。}对 $K$ 作一个\Term{finite-extension}{有限扩张}后，可以假定 $G_\eta$ 有“Chevalley \Term{resolution}{分解}”，即它是一个\Term{abelian-variety}{阿贝尔簇} $D_\eta$ 被一个\Term{smooth}{光滑}\Term{connected}{连通}\Term{linear-algebraic-group}{线性代数群} $F_\eta$ 的\Term{group-extension}{扩张}（Bourbaki 讨论班，1956/57，第 145 讲）。根据第 VI$_B$ 讲 10.16，存在 $\eta$ 的一个\Term{open-neighborhood}{开邻域} $U$，使这个\Term{generic-point}{泛点}处的\Term{group-extension}{扩张}来自一个\Term{group-extension}{扩张}
\[1\longrightarrow F\longrightarrow G|_U\longrightarrow D\longrightarrow1.\]
还可以假定 $F,D$ 在 $U$ 上\Term{smooth}{光滑}、\Term{fiber}{纤维}\Term{connected}{连通}，$F$ 在 $U$ 上\Term{affine}{仿射}，且 $D$ 在 $U$ 上\Term{proper}{固有}（EGA IV 8、9、17）。于是，对 $U$ 的每一点 $s$，$(D_s,F_s)$ 就是 $G_s$ 的“Chevalley \Term{resolution}{分解}”。因此，$G_s$ 是\Term{abelian-variety}{阿贝尔簇}（相应地，\Term{affine}{仿射}），当且仅当 $F_s$（相应地，$D_s$）是\Term{unit-group}{单位群}，而这是一个\Term{constructible}{可构造}性质（EGA IV 9.2.6.1）。

为证明 i) 的最后两个断言，根据以上结果，可以假定 $G$ 在 $S$ 上\Term{affine}{仿射}。令 $q$ 为在 $S$ 上\Term{invertible}{可逆}的素数，${}_qG$ 为 $G$ 中 $q$ 次幂映射的“核”。由\Term{affine}{仿射}\Term{algebraic-group}{代数群}的结构，容易得到 $G_s$ 是\Term{torus}{环面}（相应地，\Term{unipotent}{幺幂}），当且仅当 a) ${}_qG_s$ \Term{quasi-finite}{拟有限}，这是一个\Term{constructible}{可构造}性质（EGA IV 9.3.2）；b) ${}_qG_s$ 有 $r^q$ 个\Term{geometric-point}{几何点}，其中 $r$ 为 $G$ 对 $S$ 的\Term{relative-dimension}{相对维数}（相应地，${}_qG_s$ 只有一个点）。而函数 $s\mapsto({}_qG_s\text{ 的几何点个数})$ 是\Term{constructible}{可构造}的（EGA IV 9.7.9）。i) 的证明至此完成。

\emph{iii) 的证明。a) \Term{torus}{环面}的\Term{centralizer}{中心化子}的情形。}假设 $H_\eta=\operatorname{Centr}_{G_\eta}(T_\eta)$，其中 $T_\eta$ 为 $G_\eta$ 的\Term{torus}{环面}，来证明：对 $\eta$ 的某个\Term{neighborhood}{邻域}中的点 $s$，$H_s$ 是 $G_s$ 的某个\Term{subtorus}{子环面}的\Term{centralizer}{中心化子}。根据 i)，缩小 $S$ 后，可以
\SourcePage{295}
假定 $T_\eta$ 来自 $G$ 的一个\Term{subtorus}{子环面} $T$。此时，$Z=\operatorname{Centr}_G(T)$ 可由 $G$ 的一个\Term{group-subprescheme}{子群预概形}表示（第 XI 讲 6.11）。由于 $H,Z$ 在\Term{generic-point}{泛点}处相同，它们在 $\eta$ 的一个\Term{neighborhood}{邻域}上相同。这样就证明：使 $H_s$ 为 $G_s$ 的某个\Term{subtorus}{子环面}在 $G_s$ 中的\Term{centralizer}{中心化子}的那些点 $s$ 所成的\Term{set}{集合} $E$ 是\Term{ind-constructible}{归纳可构造}的（ind-constructible，EGA IV 1）；这一结果足以在下文\Term{lemma}{引理} 6.6 中证明 $E$ 是 $S$ 的\Term{open-subset}{开子集}，因而 $E$ 自然也是 $S$ 的\Term{locally-constructible}{局部可构造}\Term{subset}{子集}。

\emph{iii) b) \Term{cartan-subgroup}{Cartan 子群}的情形。}假设 $H_\eta$ 是 $G_\eta$ 的\Term{cartan-subgroup}{Cartan 子群}，来证明：对于 $\eta$ 的某个\Term{neighborhood}{邻域}的每一点 $s$，$H_s$ 都是 $G_s$ 的\Term{cartan-subgroup}{Cartan 子群}。群 $H_\eta$ 是 $G_\eta$ 的某个\Term{torus}{环面}在 $G_\eta$ 中的\Term{centralizer}{中心化子}，并且\Term{nilpotent}{幂零}（第 XII 讲 6.6）。由 a) 及第 VI$_B$ 讲 8.4，对于 $\eta$ 的某个\Term{neighborhood}{邻域} $U$ 的每一点 $s$，$H_s$ 都具有这些性质。所以，对 $U$ 的每一点 $s$，群 $H_s$ 与 $G_s$ 的\Term{reductive-rank}{约化秩}相同，且其唯一的\Term{maximal-torus}{极大环面}是\Term{central}{中心}的（第 XII 讲 6.7）；因此它是 $G_s$ 的一个\Term{maximal-torus}{极大环面}的\Term{centralizer}{中心化子}，即 $G_s$ 的一个\Term{cartan-subgroup}{Cartan 子群}。

现假设 $H$ 不是 $G$ 的\Term{cartan-subgroup}{Cartan 子群}，来证明：对于 $\eta$ 的某个\Term{neighborhood}{邻域} $U$ 中的点 $s$，$H_s$ 都不是 $G_s$ 的\Term{cartan-subgroup}{Cartan 子群}。鉴于上面 a) 中暂时承认的断言，只须考虑 $H$ 是一个\Term{subtorus}{子环面} $T$ 在 $G$ 中的\Term{centralizer}{中心化子}的情形。此时，$H_\eta$ 包含 $H_\eta$ 的一个\Term{cartan-subgroup}{Cartan 子群} $C_\eta$。我们刚刚看到，缩小 $S$ 后，$C_\eta$ 延拓为 $G$ 的一个\Term{cartan-subgroup}{Cartan 子群} $C$，且可以假定 $C$ 包含于 $H$。根据假设，$H_\eta$ 严格包含 $C_\eta$，所以对于 $\eta$ 的某个\Term{neighborhood}{邻域} $U$ 中的点 $s$，$H_s$ 严格包含 $C_s$（EGA IV 9.5.2）；特别地，对 $s\in U$，$H_s$ 不是 $G_s$ 的\Term{cartan-subgroup}{Cartan 子群}。

\emph{ii) 的证明。}假设 $H_\eta$ 是 $G_\eta$ 的\Term{maximal-torus}{极大环面}，令 $C_\eta$ 为其在 $G_\eta$ 中的\Term{centralizer}{中心化子}。由 i) 及 iii)，$H$ 在 $\eta$ 的某个\Term{neighborhood}{邻域} $U$ 上是\Term{torus}{环面}，且 $C=\operatorname{Centr}_{G|_U}(H|_U)$ 是 $G|_U$ 的\Term{cartan-subgroup}{Cartan 子群}。为证明 $H$ 在 $\eta$ 的某个\Term{neighborhood}{邻域}上是 $G$ 的\Term{maximal-torus}{极大环面}，可以先用 $C$ 替换 $G$，再用 $C$ 的一个 Chevalley \Term{resolution}{分解}中的线性部分 $F$ 替换它（参见 i)）。令 $q$ 为在 $S$ 上\Term{invertible}{可逆}的整数，${}_qF$ 为 $F$ 中 $q$ 次幂映射的核。由于 $F_s$ \Term{affine}{仿射}、\Term{nilpotent}{幂零}、\Term{smooth}{光滑}且\Term{connected}{连通}，$F_s$ 是其\Term{maximal-torus}{极大环面} $T_s$ 与一个\Term{unipotent-group}{幺幂群}的\Term{direct-product}{直积}（BIBLE 6-04），故 ${}_qF_s={}_qT_s$。由于 $H_\eta$ 是\Term{maximal-torus}{极大环面}，${}_qH_\eta={}_qF_\eta$，从而在 $\eta$ 的某个\Term{neighborhood}{邻域} $V$ 上有 ${}_qH={}_qF$。对 $V$ 的每一点 $s$，${}_qH_s={}_qT_s$，故 $H_s=T_s$ 是\Term{maximal-torus}{极大环面}。

现假设 $H_\eta$ 不是 $G_\eta$ 的\Term{maximal-torus}{极大环面}。由 i)，只须考虑 $H_\eta$ 是\Term{torus}{环面}的情形，再假定它包含于一个严格更大的\Term{torus}{环面} $T_\eta$。后者在 $\eta$ 的一个\Term{neighborhood}{邻域} $U$ 上延拓为一个严格包含 $H$ 的\Term{torus}{环面} $T$。特别地，对 $s\in U$，$H_s$ 不是\Term{maximal-torus}{极大环面}。

\emph{iv) 的证明。}缩小 $S$ 后，可以假定 $G$ 的\Term{center}{中心} $Z$ \Term{representability}{可表}（第 VI$_B$ 讲 10.11）且在 $S$ 上\Term{flat}{平坦}，商 $G/Z$ 也如此（同上）。性质“$H_s$ 包含 $Z_s$”是\Term{constructible}{可构造}的（EGA IV 9.5.2），且 $G_s$ 的任意\Term{parabolic-subgroup}{抛物子群}都包含 $Z_s$（第 XIV 讲 4.9 a)）；因此可以用 $G/Z$ 替换 $G$，从而假定 $G$ 在 $S$ 上\Term{affine}{仿射}（第 XII 讲 6.1 及 i)）。还可以
\SourcePage{296}
假定 $G/H$ \Term{representability}{可表}。此时，$H_s$ 是 $G_s$ 的\Term{parabolic-subgroup}{抛物子群}，当且仅当 $(G/H)_s$ \Term{proper}{固有}（BIBLE，6，\Term{theorem}{定理} 4 b)），这是一个\Term{ind-constructible}{归纳可构造}性质（EGA IV 9.3.5）。因此，$E$ 是\Term{ind-constructible}{归纳可构造}的；这一点足以证明 $E$ 是\Term{open-subset}{开集}（\Term{lemma}{引理} 6.6），从而\Term{locally-constructible}{局部可构造}。

现在考察\Term{borel-subgroup}{Borel 子群}的情形。若 $H_\eta$ 是 $G_\eta$ 的\Term{borel-subgroup}{Borel 子群}，即 $G_\eta$ 的\Term{solvable}{可解}\Term{parabolic-subgroup}{抛物子群}，则由以上结果及第 VI$_B$ 讲 8.4，这些性质在 $\eta$ 的某个\Term{neighborhood}{邻域}的每一点 $s$ 处仍成立。若 $H_\eta$ 不是 $G_\eta$ 的\Term{borel-subgroup}{Borel 子群}，为证明在 $\eta$ 的某个\Term{neighborhood}{邻域}的点 $s$ 处仍如此，根据以上结果，只须考虑 $H_\eta$ 是\Term{parabolic-subgroup}{抛物子群}的情形，再假定 $H_\eta$ 包含一个\Term{borel-subgroup}{Borel 子群} $B_\eta$。刚才已经证明，后者在 $\eta$ 的一个\Term{neighborhood}{邻域} $U$ 上延拓为 $G$ 的\Term{borel-subgroup}{Borel 子群} $B$。由于 $H_\eta$ 严格包含 $B_\eta$，在某个\Term{open-subset}{开集} $V$ 的每一点处，$H_s$ 严格包含 $B_s$，而对 $s\in V$，$H_s$ 不是 $G_s$ 的\Term{borel-subgroup}{Borel 子群}。

\emph{v) 的证明。}假设 $H_\eta$ 是 $G_\eta$ 的\Term{radical}{根基}。于是 $H_\eta$ 在 $G_\eta$ 中\Term{normal-subgroup}{正规}，且\Term{solvable}{可解}（\Term{smooth}{光滑}且\Term{connected}{连通}）；因此，对 $\eta$ 的某个\Term{neighborhood}{邻域} $U$ 中的点 $s$，$H_s$ 也如此（第 VI$_B$ 讲第 8、10 节）。所以对 $s\in U$，$H_s$ 包含于 $G_s$ 的\Term{radical}{根基}。用 $G/H$ 替换 $G$（第 VI$_B$ 讲第 10 节），只须证明：若 $G_\eta$ 是\Term{semisimple}{半单}的，则对 $\eta$ 的某个\Term{neighborhood}{邻域} $V$ 的每一点，$G_s$ 都是\Term{semisimple}{半单}的。由 i)、ii)，可以假定 $G$ 在 $S$ 上\Term{affine}{仿射}，且有一个\Term{maximal-torus}{极大环面} $T$。令 $W$ 为 $T$ 的\Term{weyl-group}{Weyl 群}（第 XII 讲 2），它在 $S$ 上\Term{quasi-finite}{拟有限}且\Term{etale}{平展}，故在某个\Term{open-subset}{开集} $V$ 上\Term{finite}{有限}且\Term{etale}{平展}。由\Term{root}{根}的初等性质（第 XIX 讲 1.12），\Term{pushout}{推出} $G$ 在 $V$ 上是\Term{semisimple}{半单}的。

现假设 $H_\eta$ 不是 $G_\eta$ 的\Term{radical}{根基}。用一个\Term{finite-extension}{有限扩张} $L$ 替换 $K$ 后，可以假定 $G_\eta$ 有\Term{radical}{根基} $R_\eta$。根据以上结果，$R_\eta$ 延拓为 $G|_U$ 的一个\Term{group-subprescheme}{子群预概形} $R$，使得对每个 $s\in U$，$R_s$ 为 $G_s$ 的\Term{radical}{根基}。由假设 $R_\eta\ne H_\eta$，故对 $s\in V$，$R_s\ne H_s$。还须证明：若 $G_\eta$ 不是\Term{semisimple}{半单}的，则对附近的点 $s$，$G_s$ 也不是\Term{semisimple}{半单}的；但这只是以上结果的一个特例（取 $H$ 为\Term{unit-group}{单位群}）。

\emph{vi) 的证明。}结合 i)，证明与 v) 完全类似，留给读者完成。
\begin{SGAresult}{推论}{6.3}
设 $S_0$ 是\Term{quasi-compact}{拟紧}\Term{prescheme}{预概形}，$S_i$（$i\in L$）是一个 $S_0$-\Term{prescheme}{预概形}的\Term{projective-system}{投射系统}，各 $S_i$ 均在 $S_0$ 上\Term{affine}{仿射}；令 $S=\varprojlim S_i$（EGA IV 8.2）。设 $G_0$ 是 $S_0$ 上\Term{finite-presentation}{有限呈示}的\Term{group-prescheme}{群预概形}，$G_i=G_0\times_{S_0}S_i$，$G=G_0\times_{S_0}S$，$H$ 是 $G$ 的一个\Term{subgroup}{子群}。若 $H$ 是 $G$ 的\Term{subtorus}{子环面}（相应地，\Term{maximal-torus}{极大子环面}、\Term{cartan-subgroup}{Cartan 子群}、\Term{borel-subgroup}{Borel 子群}、\Term{parabolic-subgroup}{抛物子群}），则存在指标 $i\in L$ 及 $G_i$ 的一个\Term{subgroup}{子群} $H_i$，使 $H=H_i\times_{S_i}S$，且 $H_i$ 为 $G_i$ 的\Term{subtorus}{子环面}（相应地，其余各情形）。
\end{SGAresult}
事实上，$H$ \Term{smooth}{光滑}且\Term{fiber}{纤维}\Term{connected}{连通}，故在 $S$ 上\Term{finite-presentation}{有限呈示}（第 VI$_B$ 讲 5.3.3）。由第 VI$_B$ 讲第 10 节，存在 $i\in L$ 及 $G_i$ 的一个在 $S$ 上\Term{smooth}{光滑}的\Term{subgroup}{子群} $H_i$，使 $H=H_i\times_{S_i}S$。推论 6.3 随即由定义 6.1、6.2 及 EGA IV 9.3.3 得到。
\SourcePage{297}
\begin{SGAresult}{推论}{6.3 bis}
设 $S$ 是\Term{prescheme}{预概形}，$G$ 是 $S$ 上\Term{finite-presentation}{有限呈示}的 $S$-\Term{group-prescheme}{群预概形}。则函数 $\rho_{\mathrm n},\rho_{\mathrm r},\rho_{\mathrm u},\rho_{\mathrm{ab}}$（参见 6.1 ter）都是 $S$ 上的\Term{locally-constructible}{局部可构造}函数。
\end{SGAresult}
只须证明（EGA IV 9）：若 $S$ 是\Term{generic-point}{泛点}为 $\eta$ 的 noether \Term{integral}{整}\Term{scheme}{概形}，则这些函数在 $\eta$ 的某个\Term{neighborhood}{邻域}上为常数。用一个在 $S$ 上\Term{finite}{有限}且\Term{dominant}{支配} $S$ 的\Term{scheme}{概形} $S'$ 替换 $S$ 后，可以假定 $G_\eta$ 有一个\Term{cartan-subgroup}{Cartan 子群} $C_\eta$，且它有一个 Chevalley \Term{resolution}{分解} $1\to L_\eta\to C_\eta\to A_\eta\to1$。6.2 i) 中的论证证明，该\Term{resolution}{分解}在 $\eta$ 的一个\Term{neighborhood}{邻域}上延拓为 Chevalley \Term{resolution}{分解}
\[1\longrightarrow L\longrightarrow C\longrightarrow A\longrightarrow1.\]
此外，可以假定 $C$ 是 $G$ 的\Term{cartan-subgroup}{Cartan 子群}（6.3），且 $L_\eta$ 的\Term{maximal-torus}{极大环面} $T_\eta$ 延拓为 $L$ 的\Term{maximal-torus}{极大环面} $T$（6.3）。由此及定义立即得到推论。

\textbf{6.4.0。}现设 $S$ 是\Term{prescheme}{预概形}，$G$ 是 $S$ 上\Term{finite-presentation}{有限呈示}的 $S$-\Term{group-prescheme}{群预概形}，$\mathscr L$ 是由 $G$ 的\Term{smooth}{光滑}、\Term{fiber}{纤维}\Term{connected}{连通}且等于自身\Term{connected-normalizer}{连通正规化子}的\Term{subgroup}{子群}所成的\Term{functor}{函子}（参见第 5 节）；$\mathscr L$ \Term{representability}{可表}（5.2、5.8）。本节余下部分研究 $\mathscr L$ 的某些\Term{subfunctor}{子函子}。具体地，引入 $\mathscr L$ 的\Term{subfunctor}{子函子} $\mathscr L\mathscr C$（相应地，$\mathscr C\mathscr T$、$\mathscr P$），定义如下：对于每个 $S$-\Term{prescheme}{预概形} $S'$，$\mathscr L\mathscr C(S')$（相应地，$\mathscr C\mathscr T(S')$、$\mathscr P(S')$）是由 $G_{S'}$ 中那些在 $S'$ 上\Term{smooth}{光滑}、\Term{fiber}{纤维}\Term{connected}{连通}、等于自身\Term{connected-normalizer}{连通正规化子}并满足下列条件的\Term{subgroup}{子群} $H$ 所成的\Term{set}{集合}：对 $S'$ 的每一点 $s$，$H_{\bar s}$ 包含 $G_{\bar s}$ 的一个\Term{cartan-subgroup}{Cartan 子群}（6.1）（相应地，$H_{\bar s}$ 是 $G_{\bar s}$ 的一个\Term{subtorus}{子环面}在 $(G_{\bar s})_{\mathrm{red}}^0$ 中的\Term{centralizer}{中心化子}；相应地，是 $G_{\bar s}$ 的一个\Term{parabolic-subgroup}{抛物子群}（6.1））。
\begin{SGAresult}{\Term{theorem}{定理}}{6.4}
设 $S$ 是\Term{prescheme}{预概形}，$G$ 是 $S$ 上\Term{finite-presentation}{有限呈示}的 $S$-\Term{group-prescheme}{群预概形}，则上述（6.4.0）$S$-\Term{functor}{函子} $\mathscr L\mathscr C,\mathscr C\mathscr T,\mathscr P$ 均可由在 $S$ 上\Term{finite-presentation}{有限呈示}且\Term{quasi-projective}{拟射影}的 $S$-\Term{prescheme}{预概形}表示。
\end{SGAresult}
\begin{SGAresult}{注记}{6.5}
若 $G$ 的\Term{fiber}{纤维}\Term{smooth}{光滑}，例如 $G$ 在 $S$ 上\Term{smooth}{光滑}，或 $S$ 的\Term{residue-characteristic}{剩余特征}均为零（第 VI$_B$ 讲 1.6.1），则 $G$ 中任一在 $S$ 上\Term{smooth}{光滑}、\Term{fiber}{纤维}\Term{connected}{连通}，且对 $S$ 的每一点 $s$ 都满足 $H_{\bar s}$ 包含 $(G_{\bar s})^0$ 的一个\Term{cartan-subgroup}{Cartan 子群}的\Term{subgroup}{子群} $H$，必然等于自身的\Term{connected-normalizer}{连通正规化子}（从而是 $\mathscr L\mathscr C(S)$ 的一个元素）。事实上，根据 5.1 iii)，可以假定 $S$ 为域的谱；在该情形中，此性质已在第 XIII 讲 2.1 陈述的末尾指出。
\end{SGAresult}
注意有自然\Term{monomorphism}{单态射}
\[
\begin{tikzcd}
\mathscr C\mathscr T\arrow[dr]&&\\
&\mathscr L\mathscr C\arrow[r]&\mathscr L.\\
\mathscr P\arrow[ur]&&
\end{tikzcd}
\]
\SourcePage{298}
下面证明这些\Term{monomorphism}{单态射}都是\Term{finite-presentation}{有限呈示}\Term{open-immersion}{开浸入}（结合 5.2，这已经说明 $\mathscr L\mathscr C,\mathscr C\mathscr T,\mathscr P$ 可由 $S$-\Term{prescheme}{预概形}表示，且它们是一列递增的\Term{open-subprescheme}{开子预概形}的并，各项均在 $S$ 上\Term{quasi-projective}{拟射影}且\Term{finite-presentation}{有限呈示}）。用通常的技巧，这将由下列\Term{lemma}{引理}得到：
\begin{SGAresult}{\Term{lemma}{引理}}{6.6}
设 $S$ 是\Term{prescheme}{预概形}，$G$ 是\Term{finite-presentation}{有限呈示} $S$-\Term{group-prescheme}{群预概形}，$H$ 是 $G$ 的一个\Term{smooth}{光滑}且\Term{fiber}{纤维}\Term{connected}{连通}的\Term{subgroup}{子群}。则 $S$ 中使 $H_{\bar s}$ 包含 $G_{\bar s}$ 的一个\Term{cartan-subgroup}{Cartan 子群}（相应地，使 $H_{\bar s}$ 等于 $G_{\bar s}$ 的某个\Term{torus}{环面}在 $(G_{\bar s})_{\mathrm{red}}^0$ 中的\Term{centralizer}{中心化子}）的点 $s$ 所成的\Term{set}{集合}是 $S$ 的\Term{open-subset}{开集}。如果还假设 $H$ 等于自身的\Term{connected-normalizer}{连通正规化子}\footnote{这个假设实际上是多余的，参见第 XVII 讲附录 III，3。}，则 $S$ 中使 $H_{\bar s}$ 为 $G_{\bar s}$ 的\Term{parabolic-subgroup}{抛物子群}的点 $s$ 所成的\Term{set}{集合}也是 $S$ 的\Term{open-subset}{开集}。
\end{SGAresult}
要证明的断言在 $S$ 上具有\Term{local-character}{局部性}，因此可以假定 $S$ \Term{affine}{仿射}，再由 EGA IV 8.1 及 6.3，假定 $S$ 为 noether。以 $E$ 表示 $S$ 中具有所讨论性质的点的\Term{set}{集合}。由 6.2 中已经实际证明的断言，$E$ 是\Term{ind-constructible}{归纳可构造}的。但 $S$ 为 noether，所以为证明 $E$ 是\Term{open-subset}{开集}，只须证明 $E$ 在\Term{generization}{一般化}下\Term{stable}{稳定}（EGA IV 1.10.1）。利用 EGA II 7.1.9，最终只须证明：若 $S$ 是\Term{discrete-valuation-ring}{离散赋值环}的谱，且\Term{closed-point}{闭点} $s$ 属于 $E$，则\Term{generic-point}{泛点} $t$ 也属于 $E$。我们需要下面的\Term{lemma}{引理}：
\begin{SGAresult}{\Term{lemma}{引理}}{6.7}
设 $A$ 是\Term{complete}{完备} noether \Term{local-ring}{局部环}，$S=\operatorname{Spec}A$，$s$ 为 $S$ 的\Term{closed-point}{闭点}，$H$ 是\Term{smooth}{光滑}且\Term{fiber}{纤维}\Term{connected}{连通}的 $S$-\Term{group-prescheme}{群预概形}，$T_s$ 是 $H_s$ 的一个\Term{subtorus}{子环面}。则：

i) 存在 $H$ 的一个\Term{smooth}{光滑}且\Term{fiber}{纤维}\Term{connected}{连通}的\Term{closed-group-subprescheme}{闭子群预概形} $C$，使 $C_s=\operatorname{Centr}_{H_s}(T_s)$。

ii) 对 $S$ 的每一点 $t$，$C_{\bar t}$ 是 $H_{\bar t}$ 的某个\Term{subtorus}{子环面} $T_{\bar t}$ 在 $H_{\bar t}$ 中的\Term{centralizer}{中心化子}。
\end{SGAresult}
\textbf{6.7 的证明。}取 $S$-\Term{torus}{环面} $T'$，使存在\Term{isomorphism}{同构} $u_0:T'_s\xrightarrow{\sim}T_s$（第 X 讲 4.6）。令 $\mathfrak m$ 为 $A$ 的\Term{maximal-ideal}{极大理想}，$A_n=A/\mathfrak m^n$，$S_n=\operatorname{Spec}A_n$，$H_n=H\times_SS_n$，其余类似。由于 $H$ 在 $S$ 上\Term{smooth}{光滑}，对每个整数 $n>0$，存在一个 $S_n$-\Term{group-morphism}{群态射}
\[u_n:T'_n\longrightarrow H_n\]
\Term{lifting}{提升} $u_0$（第 IX 讲 3.6）；对 $n$ 归纳，可以假定 $u_n$ \Term{lifting}{提升} $u_{n-1}$。另一方面，令 $q$ 为在 $S$ 上\Term{invertible}{可逆}的素数；对每个为 $q$ 的幂的整数 $\ell$，以 ${}_{\ell}u_n$ 表示 $u_n$ 在\Term{subgroup}{子群} ${}_{\ell}T'_n$ 上的\Term{restriction}{限制}。固定 $\ell$ 而令 $n$ 变化时，\Term{morphism}{态射} ${}_{\ell}u_n$ 构成\Term{projective-system}{投射系统}，故来自唯一的 $S$-\Term{group-morphism}{群态射} ${}_{\ell}u:{}_{\ell}T'\to H$（1.6 a)）。由于 $H$ \Term{separated}{分离}（第 VI$_B$ 讲 5.2），且 $u_0$ 为\Term{monomorphism}{单态射}，${}_{\ell}u$ 是\Term{monomorphism}{单态射}（第 IX 讲 6.8）；而且它是\Term{closed-immersion}{闭浸入}，因为 ${}_{\ell}T'$ 在 $S$ 上\Term{finite}{有限}，从而该\Term{morphism}{态射}\Term{finite}{有限}。以 $M(\ell)$ 表示像群。显然，\Term{multiplicative-type}{乘法型}\Term{subgroup}{子群}族 $M(\ell)$ 在 4.1 的意义下相容。令 $C_{\ell}=\operatorname{Centr}_H(M(\ell))$，它可由一个\Term{group-subprescheme}{子群预概形}表示（2.5），该\Term{subprescheme}{子预概形}闭（因为 $H$ \Term{separated}{分离}）且在 $S$ 上\Term{smooth}{光滑}（第 XI 讲 2.4）。各 $C_{\ell}$
\SourcePage{299}
构成一族递减的\Term{filtered}{滤过}\Term{closed-subscheme}{闭子概形}；由于 $H$ 为 noether，该族\Term{stable}{稳定}。其\Term{stable}{稳定}值为一个\Term{smooth}{光滑}\Term{closed-subgroup}{闭子群} $C$；由\Term{dense}{稠密}性\Term{theorem}{定理}（第 IX 讲 4.7），$C_s=\operatorname{Centr}_{H_s}(T_s)$。还须证明：对 $S$ 的每一点 $t$，$C_{\bar t}$ 为 $H_{\bar t}$ 的某个\Term{subtorus}{子环面} $T_{\bar t}$ 在 $H_{\bar t}$ 中的\Term{centralizer}{中心化子}（这将\Term{pushout}{推出} $C_{\bar t}$ \Term{connected}{连通}）。把下面更精确的\Term{lemma}{引理}应用于 $H_{\bar t}$ 的\Term{subgroup}{子群}族 $M(\ell)_{\bar t}$，即可得到此结论：
\begin{SGAresult}{\Term{lemma}{引理}}{6.8}
设 $G$ 是域 $k$ 上的\Term{connected}{连通}\Term{algebraic-group}{代数群}，$r>0$ 为整数，$q$ 为与 $k$ 的\Term{characteristic}{特征}\footnote{编者注：删去了“剩余”一词。}互素的整数，$M(\ell)$（$\ell$ 遍历 $q$ 的幂）是 $G$ 的一族相容的\Term{multiplicative-type}{乘法型}\Term{subgroup}{子群}，其型为 $(\mathbb Z/\ell\mathbb Z)^r$（参见 4.6）；$M$ 是由各 $M(\ell)$ 生成的 $G$ 的\Term{algebraic-subgroup}{代数子群}（同上），$T$ 是 $M$ 的唯一\Term{maximal-torus}{极大环面}（参见 3.4）。则当 $\ell$ \Term{large-relative-to-cardinal}{足够大}时，
\[\operatorname{Centr}_G(T)=\operatorname{Centr}_G(M)=\operatorname{Centr}_G(M(\ell)).\]
\end{SGAresult}
最后一个等号显然成立。为证明第一个，引入 $G$ 的\Term{center}{中心} $Z$、$G'=G/Z$、$M$（相应地 $T$）在 $G'$ 中的像 $M'$（相应地 $T'$），以及 $T'$ 在 $G$ 中的\Term{inverse-image}{逆像} $K$（即由 $T,Z$ 生成的 $G$ 的\Term{algebraic-subgroup}{代数子群}）。显然，只须证明 $K\supset M$，因而只须证明 $T'=M'$。而 $M$ \Term{smooth}{光滑}且\Term{connected}{连通}（4.6），$G'$ \Term{affine}{仿射}（第 XII 讲 6.1），所以 $M'$ 是其\Term{maximal-torus}{极大环面} $T'$（第 XII 讲 6.6 d)）与一个\Term{unipotent-group}{幺幂群}的\Term{direct-product}{直积}（BIBLE，4，\Term{theorem}{定理} 4）（可以假定 $k$ \Term{algebraically-closed}{代数闭}）。因此 $M(\ell)$ 在 $M'$ 中的像必然包含于 $T'$。所以，$T'$ 在 $M$ 中的\Term{inverse-image}{逆像}对所有 $\ell$ 都包含 $M(\ell)$，从而等于 $M$。因此 $M'=T'$。这就证明了 6.8，从而证明了 6.7。

现在证明 6.6。已归结到 $S$ 为\Term{discrete-valuation-ring}{离散赋值环} $A$ 的谱的情形，还可以假定 $A$ \Term{complete}{完备}，且\Term{residue-field}{剩余域}\Term{algebraically-closed}{代数闭}。用 $A$ 在其\Term{field-of-fractions}{分式域}的一个\Term{finite-extension}{有限扩张}中的\Term{normalization}{整闭包}替换 $A$，可以假定 $(G_t)_{\mathrm{red}}$ \Term{smooth}{光滑}\footnote{编者注：此处须给出细节或参考文献……\SGAsourcegap{编者明示未补的论证细节或参考文献，并非物理缺页或扫描漏页；原书此处只留下待补说明。}}。显然，为证明 6.6，可以用 $(G_t)_{\mathrm{red}}^0$ 在 $G$ 中的\Term{schematic-closure}{概形闭包}的\Term{identity-component}{单位连通分支}替换 $G$，从而假定 $G$ 在 $S$ 上\Term{flat}{平坦}、\Term{fiber}{纤维}\Term{connected}{连通}，且 $G_t$ \Term{smooth}{光滑}。

\textbf{a)} 假设 $H_{\bar s}$ 是一个\Term{torus}{环面} $T_{\bar s}$ 在 $(G_{\bar s})_{\mathrm{red}}$ 中的\Term{centralizer}{中心化子}，来证明 $H_{\bar t}$ 是 $G_{\bar t}$ 的某个\Term{subtorus}{子环面}在 $G_{\bar t}$ 中的\Term{centralizer}{中心化子}。由\Term{lemma}{引理} 6.7，存在 $H$ 的一个在 $S$ 上\Term{smooth}{光滑}的\Term{group-subscheme}{子群概形} $C$，其闭\Term{fiber}{纤维}为 $H_s$，并使 $C_{\bar t}=\operatorname{Centr}_{H_{\bar t}}(T_{\bar t})$，其中 $T_{\bar t}$ 是 $H_{\bar t}$ 的一个\Term{subtorus}{子环面}。由于 $H$ 在 $S$ 上\Term{smooth}{光滑}且\Term{fiber}{纤维}\Term{connected}{连通}，由\Term{dimension}{维数}的考量可得 $C=H$。沿用 6.7 的记号，当 $\ell$ \Term{large-relative-to-cardinal}{足够大}时，有 $H=C=C_{\ell}$。同样，考察 $C'_{\ell}=\operatorname{Centr}_G(M(\ell))$（2.5），令 $C'$ 为 $\ell$ \Term{large-relative-to-cardinal}{足够大}时 $C'_{\ell}$ 的\Term{stable}{稳定}值（2.5 bis）。\Term{group-scheme}{群概形} $C'$ 包含 $H$，并且 $C'_{\bar t}=\operatorname{Centr}_{G_{\bar t}}(T_{\bar t})$（6.8），$C'_{\bar s}=\operatorname{Centr}_{G_{\bar s}}(T_{\bar s})$。对 $H_s$ 的假设给出 $\dim H_s=\dim C'_s$。另一方面，$\dim H_s=\dim H_t$（因为 $H$ 在 $S$ 上\Term{smooth}{光滑}），且 $\dim C'_t\leq\dim C'_s$（第 VI$_B$ 讲 4.1），故 $\dim H_t=\dim C'_t$。但由于 $G_{\bar t}$ \Term{smooth}{光滑}且\Term{connected}{连通}，$C'_{\bar t}=\operatorname{Centr}_{G_{\bar t}}(T_{\bar t})$ \Term{smooth}{光滑}且\Term{connected}{连通}，所以最终有 $H_{\bar t}=C'_{\bar t}=\operatorname{Centr}_{G_{\bar t}}(T_{\bar t})$。
\SourcePage{300}
\textbf{b)} 假设 $H_{\bar s}$ 包含 $G_{\bar s}$ 的一个\Term{cartan-subgroup}{Cartan 子群} $C_{\bar s}$，即 $G_{\bar s}$ 的一个\Term{maximal-torus}{极大环面} $T_{\bar s}$ 在 $(G_{\bar s})_{\mathrm{red}}$ 中的\Term{centralizer}{中心化子}。由 6.7，存在 $H$ 的一个在 $S$ 上\Term{smooth}{光滑}且\Term{fiber}{纤维}\Term{connected}{连通}的\Term{group-subscheme}{子群概形} $C$，它\Term{lifting}{提升} $C_s$。显然，只须证明 $C_{\bar t}$ 包含 $G_{\bar t}$ 的一个\Term{cartan-subgroup}{Cartan 子群}。而取 $H=C$ 应用 a)，$C_{\bar t}$ 是 $G_{\bar t}$ 的某个\Term{subtorus}{子环面}在 $G_{\bar t}$ 中的\Term{centralizer}{中心化子}，因而包含 $G_{\bar t}$ 的一个\Term{cartan-subgroup}{Cartan 子群}。

\textbf{c)} 假设 $H_{\bar s}$ 是 $G_{\bar s}$ 的\Term{parabolic-subgroup}{抛物子群}。令 $N=\operatorname{Norm}_G(H)$。由假设，$H$ 等于自身的\Term{connected-normalizer}{连通正规化子}，所以 $N_s$ \Term{smooth}{光滑}，进而等于 $\operatorname{Norm}_{(G_s)_{\mathrm{red}}}(H_s)=H_s$（第 XII 讲 8 bis）。此时 $N$ 在 $S$ 上\Term{flat}{平坦}。我们将在第 XVI 讲看到，在这些条件下，$G/N$ \Term{representability}{可表}。由于 $H_s=N_s$ 是 $G$ 的\Term{parabolic-subgroup}{抛物子群}，$(G/N)_s$ \Term{proper}{固有}。又由于 $G/N$ 的\Term{fiber}{纤维}\Term{connected}{连通}，且在 $S$ 上\Term{flat}{平坦}，由 EGA III 5.5.1，$G/N$ 在 $S$ 上\Term{proper}{固有}。因此 $(G/N)_t=G_t/N_t$ 在 $t$ 上\Term{proper}{固有}，由于 $N_t/H_t$ \Term{finite}{有限}，$G_t/H_t$ 也\Term{proper}{固有}。只须应用下面的\Term{lemma}{引理}：
\begin{SGAresult}{\Term{lemma}{引理}}{6.9}
设 $k$ 是域，$G$ 是 $k$-\Term{algebraic-group}{代数群}，$H$ 是 $G$ 的一个\Term{smooth}{光滑}\Term{connected}{连通}\Term{algebraic-subgroup}{代数子群}，$N=\operatorname{Norm}_GH$；若 $\dim H=\dim N$，且 $G/H$ \Term{proper}{固有}，则 $H$ 是 $G$ 的\Term{parabolic-subgroup}{抛物子群}。
\end{SGAresult}
事实上，可以假定 $k$ \Term{algebraically-closed}{代数闭}，$G$ \Term{smooth}{光滑}且\Term{connected}{连通}。$G$ 的\Term{center}{中心} $Z$ 包含于 $N$，而假设 $\dim H=\dim N$ \Term{pushout}{推出} $Z'=(Z)_{\mathrm{red}}^0$ 包含于 $H$，所以 $G'=G/Z'$ \Term{affine}{仿射}（第 XII 讲 6.1）。用 $G'$ 替换 $G$，用 $H$ 在 $G'$ 中的像替换 $H$，就归结到 $G$ \Term{affine}{仿射}的情形（第 XIV 讲 4.9）；此时\Term{lemma}{引理} 6.9 由 BIBLE，6，\Term{theorem}{定理} 4 得到。\Term{lemma}{引理} 6.6 因而得证。

为完成 6.4 的证明，还须证明：表示 $\mathscr L\mathscr C$（相应地，$\mathscr C\mathscr T$、$\mathscr P$）的 $S$-\Term{prescheme}{预概形}在 $S$ 上\Term{finite-presentation}{有限呈示}。这个断言在 $S$ 上具有\Term{local-character}{局部性}，因此可以假定 $S$ \Term{affine}{仿射}；再由于 $G$ \Term{finite-presentation}{有限呈示}，可以假定 $S$ 为 noether（EGA IV 8.9）。刚才已经看到，自然包含 $\mathscr C\mathscr T\to\mathscr L$（相应地，$\mathscr P\to\mathscr L$）是\Term{immersion}{浸入}，故包含 $\mathscr C\mathscr T\to\mathscr L\mathscr C$（相应地，$\mathscr P\to\mathscr L\mathscr C$）也如此。因此，只须证明 $\mathscr L\mathscr C$ 可由一个\Term{finite-presentation}{有限呈示} $S$-\Term{prescheme}{预概形}表示。

重新采用 5.2 中引入的记号。对于每个整数 $n\geq0$，令 $\mathscr L^n$ 为 $\mathscr L$ 的\Term{subfunctor}{子函子}，使
\[\mathscr L^n(S')=\{H\in\mathscr L(S')\mid\operatorname{Norm}_{G_{S'}}(H)=\operatorname{Norm}_{G_{S'}}H^{(n)}\}.\]
所以 $S$-\Term{functor}{函子} $\mathscr L^n$ 可由 $\mathscr L$ 的一个\Term{open-subprescheme}{开子预概形}表示，它是各 $\mathscr L_r^n$ 的和。每个 $\mathscr L_r^n$ 都在 $S$ 上\Term{finite-presentation}{有限呈示}（5.3），且当 $r>\sup_{s\in S}\dim G_s$ 时为空（此\Term{supremum}{上确界}\Term{finite}{有限}，因为 $S$ \Term{quasi-compact}{拟紧}），因此 $\mathscr L^n$ 在 $S$ 上\Term{finite-presentation}{有限呈示}。只须证明：当 $n$ \Term{large-relative-to-cardinal}{足够大}时，$\mathscr L\mathscr C$ 包含于 $\mathscr L^n$。

对于 $S$ 的每一点 $s$，令 $d(s)$ 为使 $\mathscr L\mathscr C_{G_s}\subset\mathscr L_{G_s}^n$ 成立的最小整数 $n$（\Term{finite}{有限}或无限）。只须证明函数 $d$ 在 $S$ \Term{bounded-above}{上有界}：因为若 $M$ 是一个上界，则 $\mathscr L\mathscr C$ 在\Term{set}{集合}意义下包含于 $\mathscr L^M$，而后者是 $\mathscr L$ 的\Term{open-subset}{开集}，所以 $\mathscr L\mathscr C$ 就包含于 $\mathscr L^M$。一个直接的\Term{constructibility}{可构造性}论证
\SourcePage{301}
使我们归结到证明：若 $S$ 是\Term{generic-point}{泛点}为 $t$ 的 noether \Term{integral}{整}\Term{prescheme}{预概形}，则函数 $d$ 在 $t$ 的某个\Term{neighborhood}{邻域}\Term{bounded-above}{上有界}。

\textbf{a) 归结到 $G$ 在 $S$ 上\Term{smooth}{光滑}的情形。}

如 6.2 那样进行，改变 $S$ 后，可以假定 $(G)_{\mathrm{red}}$ 是 $S$ 上的\Term{smooth}{光滑}\Term{group-prescheme}{群预概形}，记为 $G'$。置 $X=\mathscr L\mathscr C_G$，$X'=\mathscr L\mathscr C_{G'}$，令 $H$（相应地，$H'$）为 $G_X$（相应地，$G'_{X'}$）中对\Term{functor}{函子} $\mathscr L\mathscr C_G$（相应地，$\mathscr L\mathscr C_{G'}$）具有\Term{versal}{万有}性的\Term{group-subprescheme}{子群预概形}。由于 $H$ 在 $X$ 上\Term{smooth}{光滑}，$H\times_XX_{\mathrm{red}}$ \Term{reduced}{约化}，故包含于 $G'_{X_{\mathrm{red}}}$，且是 $\mathscr L\mathscr C_{G'}(X_{\mathrm{red}})$ 的一个元素；由此得到\Term{canonical}{典范}\Term{morphism}{态射} $p:X_{\mathrm{red}}\to X'$。显然，$p$ 是\Term{monomorphism}{单态射}；下面证明 $p$ 甚至是\Term{immersion}{浸入}。令 $N'$ 为 $H'$ 在 $G_{X'}$ 中的\Term{normalizer}{正规化子}。$X'$ 中使\Term{immersion}{浸入} $H'_s\to N'_s$ 为\Term{open-immersion}{开浸入}的点 $s$ 所成的\Term{set}{集合}是一个\Term{open-subset}{开集} $U$，且 $H'|_U\to N|_U$ 是\Term{open-immersion}{开浸入}（第 VI$_B$ 讲 2.5 及 EGA IV 17.9.5）。由 5.1 iii)，$U$ 是 $X'$ 中使 $H'$ 在其上等于自身在 $G_{X'}$ 中的\Term{connected-normalizer}{连通正规化子}的最大\Term{open-subset}{开集}，所以 $H'|_U\in\mathscr L\mathscr C_G(U)$。立即\Term{pushout}{推出} $p$ 是从 $X_{\mathrm{red}}$ 到 $U_{\mathrm{red}}$ 的\Term{isomorphism}{同构}。若能证明当 $G$ \Term{smooth}{光滑}时，$\mathscr L\mathscr C_G$ \Term{finite-type}{有限型}，则 $X'$ 在 $S$ 上\Term{finite-type}{有限型}，故 $X_{\mathrm{red}}$ 也\Term{finite-type}{有限型}（因为 $S$ 为 noether），从而当 $M$ \Term{large-relative-to-cardinal}{足够大}时包含于 $\mathscr L_G^M$，而 $X$ 也如此。

\textbf{b) $S$ 是\Term{characteristic}{特征} $p>0$ 的\Term{algebraically-closed-field}{代数闭域} $k$ 的谱，且 $G$ 是 $k$ 上\Term{smooth}{光滑}\Term{algebraic-group}{代数群}的情形。}

我们将使用 $H$ 的\Term{radicial}{根性}\Term{subgroup}{子群} ${}_{\mathrm F^n}(H)$，即 $H$ 中 Frobenius \Term{morphism}{态射}的迭代的核（第 VII$_A$ 讲 4），来代替 $G$ 的\Term{group-subprescheme}{子群预概形} $H$ 的\Term{infinitesimal-neighborhood}{无穷小邻域} $H^{(n)}$。这里这样做是允许的，因为 ${}_{\mathrm F^n}(H)$ 包含于 $H$ 的\Term{unit-section}{单位截面}的 $p^n$ 阶\Term{infinitesimal-neighborhood}{无穷小邻域}。若 $T$ 是 $G$ 的\Term{subtorus}{子环面}，则 ${}_{\mathrm F^n}(T)={}_{p^n}(T)$；由\Term{duality}{对偶性}立即可知，当 $n$ \Term{large-relative-to-cardinal}{足够大}时，$\operatorname{Centr}_G({}_{p^n}(T))$ 等于 $\operatorname{Centr}_G(T)$。于是有下列更精确的命题：
\begin{SGAresult}{命题}{6.10}
设 $k$ 是\Term{characteristic}{特征} $p>0$ 的域，$G$ 是\Term{smooth}{光滑} $k$-\Term{algebraic-group}{代数群}，$T$ 是 $G_{\bar k}$ 的一个\Term{maximal-torus}{极大环面}，$m$ 是使
\[\operatorname{Centr}_G({}_{p^m}(T))^0=\operatorname{Centr}_G(T)^0\]
成立的最小整数。则对于每个 $k$-\Term{prescheme}{预概形} $S$ 以及每个 $H\in\mathscr L\mathscr C(S)$，有
\[\operatorname{Norm}_{G_S}(H)=\operatorname{Norm}_{G_S}({}_{\mathrm F^m}(H)).\]
特别地，$\mathscr L\mathscr C$ 包含于 $\mathscr L^{p^m}$。
\end{SGAresult}
由于 $H$ 在 $S$ 上\Term{smooth}{光滑}、\Term{finite-presentation}{有限呈示}，且\Term{nilpotent-rank}{幂零秩}为常数（即 $G$ 的\Term{nilpotent-rank}{幂零秩}），我们将在下一节（7.3）看到，$H$ 的\Term{cartan-subgroup}{Cartan 子群}\Term{functor}{函子} $\mathscr C_H$ 可由一个在 $S$ 上\Term{smooth}{光滑}的 $S$-\Term{prescheme}{预概形}表示（读者可以检查，该性质的证明没有用到 $\mathscr L\mathscr C_H$ 在 $S$ 上\Term{finite-type}{有限型}这一点）。由第 XIII 讲 3.1，可以考察 $H$ 的\Term{regular-element}{正则点}\Term{open-subset}{开集} $U$。

设 $S'$ 是 $S$-\Term{prescheme}{预概形}，$g\in G(S')$ \Term{normalization}{正规化} ${}_{\mathrm F^m}(H)_{S'}$。为证明 $g$ \Term{normalization}{正规化} $H_{S'}$，只须证明 $\operatorname{int}(g)U_{S'}$ 包含于 $H_{S'}$；因为此时 $H_{S'}\cap\operatorname{int}(g)H_{S'}$ 将包含 $\operatorname{int}(g)H$ 的一个\Term{open-subgroup}{开子群}，从而等于
\SourcePage{302}
$\operatorname{int}(g)H$，因为后者的\Term{fiber}{纤维}\Term{connected}{连通}。用适当的 $S''$ 替换 $S'$，再将 $S''$ 记为 $S$，便归结为证明：若 $g\in G(S)$ \Term{normalization}{正规化} ${}_{\mathrm F^m}(H)$，且 $u\in U(S)$，则 $\operatorname{int}(g)u\in H(S)$。

令 $C$ 为 $H_S$ 中“包含”$u$ 的唯一\Term{cartan-subgroup}{Cartan 子群}（第 XIII 讲 3.2）。只须证明 $C'=\operatorname{int}(g)C\subset H$。由于 $H\in\mathscr L\mathscr C_G(S)$，$C$ 也是 $G$ 的\Term{cartan-subgroup}{Cartan 子群}；而后者具有\Term{maximal-torus}{极大环面}，所以（第 XII 讲 7.1 (a)）$C$ 是它的唯一\Term{maximal-torus}{极大环面} $T$ 在 $G_S^0$ 中的\Term{centralizer}{中心化子}。由 $m$ 的定义（以及 $G$ 的任意两个\Term{maximal-torus}{极大环面}对于 fpqc \Term{local-scheme}{局部}\Term{conjugacy}{共轭}这一事实，第 XII 讲 7.1），有
\[C=(\operatorname{Centr}_G(T))^0=\operatorname{Centr}_G({}_{p^m}(T))^0,\]
用 $g$ \Term{conjugacy}{共轭}，得到
\[C'=\operatorname{Centr}_G\bigl(\operatorname{int}(g)({}_{p^m}(T))\bigr)^0.\]
而 $\operatorname{int}(g)({}_{p^m}(T))$ 是 $\operatorname{int}(g)({}_{\mathrm F^m}(H))$ 的\Term{multiplicative-type}{乘法型}\Term{subgroup}{子群}，后者等于 ${}_{\mathrm F^m}(H)$（因为 $g$ \Term{normalization}{正规化} ${}_{\mathrm F^m}(H)$），故它包含于 $H$。由第 XIII 讲 2.1（该结果是在\Term{base}{基底}为域时证明的，但立即可推广到任意\Term{base}{基底}的情形），为使 $C'$ 包含于 $H$，只须有
\[\operatorname{Lie}C'\subset\operatorname{Lie}H.\]
而 $\operatorname{Lie}C'=\operatorname{int}(g)(\operatorname{Lie}C)\subset\operatorname{int}(g)(\operatorname{Lie}H)$。

另一方面，若 $m\geq1$（可以作此假设），则利用第 VII$_A$ 讲 4.1.2，有
\[\operatorname{Lie}H=\operatorname{Lie}({}_{\mathrm F^m}(H))=\operatorname{Lie}\bigl(\operatorname{int}(g)({}_{\mathrm F^m}(H))\bigr)=\operatorname{int}(g)(\operatorname{Lie}H).\]
因此 $\operatorname{Lie}C'\subset\operatorname{Lie}H$，6.10 的证明至此完成。

我们需要对 6.10 中引入的整数 $m$ 给出另一个定义。
\begin{SGAresult}{\Term{lemma}{引理}}{6.11}
设 $G$ 是\Term{characteristic}{特征} $p>0$ 的\Term{algebraically-closed-field}{代数闭域} $k$ 上的\Term{smooth}{光滑}\Term{algebraic-group}{代数群}，$T$ 是 $G$ 的\Term{maximal-torus}{极大环面}，$\mathfrak g=\operatorname{Lie}G$；令 $R$ 为 $T$ 的那些非零\Term{character}{特征标}所成的族，它们出现在由 $G$ 的\Term{adjoint-representation}{伴随表示}诱导的 $T$ 在 $\mathfrak g$ 上的表示中。对于每个 $r\in R$，以 $e_r$ 表示使 $p^n$ 在 $T$ 的\Term{character-group}{特征标群}中整除 $r$ 的最大整数 $n$。则当 $R\ne\varnothing$ 时，$m=\sup_{r\in R}(e_r+1)$；否则，$m=0$。
\end{SGAresult}
事实上，$\operatorname{Centr}_G(T)$ \Term{smooth}{光滑}且包含于 $\operatorname{Centr}_G({}_{p^m}(T))$，所以
\begin{align*}
(\operatorname{Centr}_GT)^0=(\operatorname{Centr}_G({}_{p^m}(T)))^0
&\Longleftrightarrow\operatorname{Lie}(\operatorname{Centr}_GT)=\operatorname{Lie}(\operatorname{Centr}_G({}_{\mathrm F^m}(T)))\\
&\Longleftrightarrow\mathfrak g^T=\mathfrak g^{{}_{p^m}(T)}\qquad\text{（第 II 讲 5.2.3）}.
\end{align*}
采用通常的记号，有
\[\mathfrak g=\mathfrak g_0\oplus\coprod_{r\in R}\mathfrak g^r.\]
因此，$\mathfrak g^T=\mathfrak g_0$，且 $\mathfrak g^{({}_{p^m}(T))}=\mathfrak g_0+\coprod_{r\in R'}\mathfrak g^r$，其中 $R'$ 是 $R$ 中由那些在 ${}_{p^m}(T)$ 上\Term{restriction}{限制}为零的 $T$ 的\Term{character}{特征标}所成的\Term{subset}{子集}。但 $T$ 的一个非零\Term{character}{特征标} $r$ 在 ${}_{p^m}(T)$ 上\Term{restriction}{限制}为零，当且仅当 $m\leq e_r$，\Term{lemma}{引理}得证。
\SourcePage{303}
\textbf{c)} 回到 6.4 的证明。根据 a) 及 a) 之前的段落，已经归结到 $S$ 是 noether \Term{integral}{整}\Term{scheme}{概形}且 $G$ 在 $S$ 上\Term{smooth}{光滑}的情形。须证明函数 $d$ 在 $S$ 的\Term{generic-point}{泛点} $t$ 的某个\Term{neighborhood}{邻域}\Term{bounded-above}{上有界}。改变 $S$ 后，可以假定 $G$ 有一个\Term{split-torus}{平凡}的\Term{maximal-torus}{极大环面} $T=\mathrm D_S(M)$（6.2）。令 $\mathfrak g=\coprod_{\lambda\in M}\mathfrak g^{\lambda}$ 为 $G$ 的\Term{lie-algebra}{Lie 代数}按 $T$ 的\Term{character}{特征标}的\Term{resolution}{分解}，$R$ 为 $M$ 中使 $\mathfrak g^{\lambda}\ne0$ 的非零\Term{character}{特征标}所成的\Term{finite}{有限}\Term{set}{集合}。分两种情形：

\emph{第一种情形：}点 $t$，因而 $S$ 的所有点，均有\Term{residue-characteristic}{剩余特征} $p>0$。由以上结果显然可知，此时函数 $d$ 以 $p^m$ 为上界，其中 $m$ 如 6.11 中定义。

\emph{第二种情形：}点 $t$ 的\Term{residue-characteristic}{剩余特征}为零。对于每个 $\lambda\in R$，令 $n_{\lambda}$ 为在群 $M$ 中整除 $\lambda$ 的最大整数，置 $n=\prod_{\lambda\in R}n_{\lambda}$。对于每个整除 $n$ 的素数 $q$，令 $S_q$ 为 $S$ 中由\Term{residue-characteristic}{剩余特征}等于 $q$ 的点所成的\Term{closed-subset}{闭子集}，$U$ 为各 $S_q$ 之并在 $S$ 中的补集；这是一个\Term{nonempty}{非空}\Term{open-subset}{开集}（它包含 $t$）。现若 $s\in U$，则或者 $s$ 的\Term{residue-characteristic}{剩余特征}为零，此时 $d(s)=1$（5.6）；或者 $s$ 的\Term{residue-characteristic}{剩余特征}为不整除 $n$ 的 $p>0$，此时 6.11 中为群 $G_s$ 定义的整数 $m$ 至多为一。此外，由第 VII 讲\footnote{编者注：须详述这一论证……\SGAsourcegap{编者明示未补的论证细节，并非物理缺页或扫描漏页；原书此处只留下论证待详述的说明。}}，有
\[\operatorname{Norm}_{G_{S'}}({}_{\mathrm F^1}(H))=\operatorname{Norm}_{G_{S'}}(\operatorname{Lie}H)=\operatorname{Norm}_{G_{S'}}(H^{(1)}).\]
最后，由 6.10，若 $H\in\mathscr L\mathscr C_{G_s}(S')$，则 $\operatorname{Norm}_{G_{S'}}(H)=\operatorname{Norm}_{G_{S'}}(H^{(1)})$；因此 $d(s)\leq1$，故 $d$ 在 $U$ 上以 $1$ 为上界。

6.4 的证明至此完成。
\begin{SGAresult}{推论}{6.12}
设 $S$ 是\Term{prescheme}{预概形}，$G$ 是\Term{finite-presentation}{有限呈示} $S$-\Term{group-prescheme}{群预概形}。假设 $G$ 的\Term{fiber}{纤维}的\Term{nilpotent-rank}{幂零秩}（相应地，\Term{borel-subgroup}{Borel 子群}的\Term{dimension}{维数}）是 $S$ 上的\Term{locally-constant}{局部常值}函数，则 $G$ 的\Term{cartan-subgroup}{Cartan 子群}\Term{functor}{函子} $\mathscr C$（相应地，$G$ 的\Term{borel-subgroup}{Borel 子群}\Term{functor}{函子} $\mathscr B$）可由一个在 $S$ 上\Term{finite-presentation}{有限呈示}的 $S$-\Term{prescheme}{预概形}表示。
\end{SGAresult}
事实上，缩小 $S$ 后，可以假定\Term{fiber}{纤维}的\Term{nilpotent-rank}{幂零秩} $\nu$ 为常数。此时显然，$\mathscr C$ 可由表示 $\mathscr L\mathscr C$ 的\Term{prescheme}{预概形} $X$（6.4）的一个既开又闭的\Term{subprescheme}{子预概形}表示，在这个\Term{subprescheme}{子预概形}上，$G_X$ 中关于\Term{functor}{函子} $\mathscr L\mathscr C$ 的\Term{versal}{万有}\Term{subgroup}{子群}\Term{relative-dimension}{相对维数}为 $\nu$。结合\Term{functor}{函子} $\mathscr P$ 的\Term{representability}{可表性}，\Term{functor}{函子} $\mathscr B$ 的证明类似。
